\chapter{Classical Rigid-Body Realization}
\label{ch:ch04_rigidbody}

The question of this chapter is whether the interior dynamics of the arena is a rigid-body system, and what an identification with
one preserves. \Cref{sec:rigidbody} recalls conventional rigid-body mechanics, against which the identification is tested.
\Cref{sec:shifted,sec:neutral,sec:closure,sec:restclosure} contain the arena's interior dynamics: the shifted Lagrangian in a positive form,
binding in neutral signature, the closure of the force law into a unique rotor, and the rest energy in the closure. \Cref{sec:esmap} tests
the identification against equations, brackets, conserved quantities and degrees of freedom, and \cref{sec:orbital} compares the interior with
planetary orbital dynamics. The closure is a Lie--Poisson system of rigid-body type, on a group of one rotation and six translations, and its
reduced motion is a pendulum (\cref{thm:esmap}). Its rotation, however, is an internal phase with no axis, its
Lorentz-invariant inertia is indefinite, and the bounded orbits of the interior are Hooke's, not Newton's (\cref{prop:hooke}). The algebra is proved and checked; every physical reading is
marked as one. Bounded orbits do not make the Hamiltonian bounded below: the interior energy has signature $(6,6)$, and the bound orbits are not
bound states (\cref{mb:noground}). The gravitational potential of \cref{sec:orbital} is assumed.

\section{Conventional rigid-body mechanics}
\label{sec:rigidbody}

This section recalls the mechanics of a rigid body in the form in which it will be compared with the arena: the inertia tensor, angular
velocity and momentum, Euler's equations as an Euler--Arnold equation, the Lie--Poisson structure and its conservation laws, and the planar
body that rotates and translates. Everything here is inherited \cite{Euler1765,Arnold66,ArnoldMMCM,LandauLifshitz1,MarsdenRatiu99}; the
checks rb1-1 to rb1-8 only fix the signs and conventions used later.

\subsection{Configuration, angular velocity and inertia}
\label{sec:rb:config}

A rigid body with one point held fixed has the configuration space $SO(3)$: a material point with body coordinates $\mathbf X$ is at
$\mathbf x(t)=R(t)\mathbf X$, $R\in SO(3)$. The body angular velocity $\boldsymbol\Omega$ is defined by $\dot R=R\,\widehat{\boldsymbol\Omega}$, where
$\widehat{\boldsymbol\Omega}\mathbf v=\boldsymbol\Omega\times\mathbf v$, and the spatial one is $\boldsymbol\omega=R\boldsymbol\Omega$. For a mass density $\rho$ the
inertia tensor is
\begin{equation}
\label{eq:rb:inertia}
\mathbb I=\int\rho(\mathbf X)\bigl(|\mathbf X|^2\,\id-\mathbf X\mathbf X^{\!\top}\bigr)\,d^3X ,
\end{equation}
symmetric and positive definite unless the body lies on a line. In principal axes its moments satisfy $I_j+I_k-I_i=2\int\rho X_i^2\,d^3X\ge0$, the
triangle inequalities (rb1-4). The kinetic energy is $T=\tfrac12\boldsymbol\Omega\cdot\mathbb I\boldsymbol\Omega$, the body angular momentum is
$\mathbf M=\mathbb I\boldsymbol\Omega$ and the spatial one is $\mathbf m=R\mathbf M$. The units are $[\mathbb I]=\mathsf{ML}^2$,
$[\boldsymbol\Omega]=\mathsf T^{-1}$ and $[\mathbf M]=\mathsf{ML}^2\mathsf T^{-1}$; the evolution parameter is the Newtonian time $t$, and the action of the
free body is $\int\tfrac12\boldsymbol\Omega\cdot\mathbb I\boldsymbol\Omega\,dt$.

\subsection{Euler's equations as an Euler--Arnold equation}
\label{sec:rb:euler}

The Lagrangian $\tfrac12\boldsymbol\Omega\cdot\mathbb I\boldsymbol\Omega$ is invariant under $R\mapsto QR$, $Q\in SO(3)$, and depends on $R$ only through
$\boldsymbol\Omega=(R^{-1}\dot R)^\vee$. Varying $R$ with $\widehat{\boldsymbol\Sigma}=R^{-1}\delta R$ gives $\delta\boldsymbol\Omega=\dot{\boldsymbol\Sigma}+\boldsymbol\Omega\times\boldsymbol\Sigma$, and
Hamilton's principle becomes Euler's equations
\begin{equation}
\label{eq:rb:euler}
\dot{\mathbf M}=\mathbf M\times\boldsymbol\Omega,\qquad\text{that is}\qquad I_1\dot\Omega_1=(I_2-I_3)\,\Omega_2\Omega_3\quad\text{and cyclically}
\end{equation}
(rb1-2). This is the case $G=SO(3)$ of Arnold's description of a left-invariant Lagrangian $\tfrac12\langle\xi,A\xi\rangle$ on a Lie group: with
$\mu=A\xi$ the motion obeys $\dot\mu=\ad_\xi^{\!\top}\mu$, and the solutions are the geodesics of the left-invariant metric $A$
\cite{Arnold66,ArnoldKhesin98}. The variations $\delta\xi=\dot\Sigma+[\xi,\Sigma]$ are those of Euler--Poincar\'e reduction \cite{MarsdenRatiu99,HolmMarsdenRatiu98}.
The steady rotations about the axes of largest and smallest moment are stable, and those about the intermediate axis are unstable: for
$(I_1,I_2,I_3)=(1,2,3)$ the linearisations have characteristic polynomials $\lambda(\lambda^2+\tfrac13\Omega^2)$, $\lambda(\lambda^2-\tfrac13\Omega^2)$ and
$\lambda(\lambda^2+\Omega^2)$ (rb1-5).

\subsection{Lie--Poisson structure and conservation laws}
\label{sec:rb:lp}

On $\so3^*\cong\R^3$ the Lie--Poisson bracket and the Hamiltonian are
\begin{equation}
\label{eq:rb:lp}
\{F,K\}(\mathbf M)=-\mathbf M\cdot(\nabla F\times\nabla K),\qquad H=\tfrac12\,\mathbf M\cdot\mathbb I^{-1}\mathbf M ,
\end{equation}
and $\dot{\mathbf M}=\{\mathbf M,H\}$ is \eqref{eq:rb:euler} (rb1-1, rb1-2). Three conservation laws follow, of three different kinds (rb1-3). The
energy $H$ is conserved because it generates the flow. The Casimir $|\mathbf M|^2$ Poisson-commutes with every function, whatever the Hamiltonian;
its level sets, the spheres, are the symplectic leaves. The spatial angular momentum $\mathbf m=R\mathbf M$ is the Noether charge of the left
action $R\mapsto QR$. Reduction thus takes the six-dimensional $T^*SO(3)$ to the three-dimensional $\so3^*$ and then to two-dimensional
leaves: one degree of freedom, and the orbits are the intersections of the spheres with the energy ellipsoids.

\subsection{Rotation with translation: the planar body}
\label{sec:rb:se2}

A body that also moves has the group of rigid motions as its configuration space. In the plane this is $SE(2)=SO(2)\ltimes\R^2$, with the
product $(\theta,\mathbf x)(\theta',\mathbf x')=(\theta+\theta',\mathbf x+e^{\theta J}\mathbf x')$, $J=\bigl(\begin{smallmatrix}0&-1\\1&0\end{smallmatrix}\bigr)$, and
the Lie algebra $\R\hat g\ltimes_J\R^2$ with $[\hat g,X]=JX$. For the Lagrangian $\tfrac12I\dot\theta^2+\tfrac12\,\mathbf u^{\!\top}N^{-1}\mathbf u$ in the
body velocity $\mathbf u=e^{-\theta J}\dot{\mathbf x}$, with $I>0$ and $N^{-1}$ a positive translational inertia, the coordinates $(h,p)$ of
$\mathfrak{se}(2)^*$ have the brackets $\{h,p_1\}=-p_2$, $\{h,p_2\}=p_1$, $\{p_1,p_2\}=0$ and the Casimir $|p|^2$ (rb1-6). With
$H=h^2/(2I)+\tfrac12p^{\!\top}Np$,
\begin{equation}
\label{eq:rb:se2}
\dot p=\frac hI\,J^{\!\top}p,\qquad \dot h=-\,p^{\!\top}JNp,\qquad \dot\theta=\frac hI,\qquad \dot{\mathbf x}=e^{\theta J}Np .
\end{equation}
On a leaf $p=e^{-\theta J}p_0$. With $N=\operatorname{diag}(a,b)$ and $p_0=(|p_0|,0)$ the rotation obeys $I\ddot\theta=\tfrac12(a-b)|p_0|^2\sin2\theta$, a
pendulum in the doubled angle; with $a=b$ the body rotates uniformly and translates uniformly, the two motions decoupled (rb1-7). The spatial
momentum $p_s=e^{\theta J}p$ and the quantity $J_s=h+p_s^{\!\top}J\mathbf x$, the total angular momentum about the origin, are conserved for every
$N$ (rb1-8). For a body symmetric about two perpendicular axes this is Kirchhoff's description of planar motion through an ideal fluid at rest at
infinity, with the mass of the body and the added masses of the fluid in $N^{-1}$ \cite{Kirchhoff1870,Lamb1932}. It is the benchmark that \cref{sec:esmap} needs: the closure of \cref{sec:closure} has the same form, with three
planes in place of one.

In this chapter a system is \emph{of rigid-body type} when it is the Lie--Poisson system of a left-invariant Lagrangian on a group of rotations,
or of rotations and translations. A \emph{physical} rigid body has a positive-definite inertia and a rotation group acting on the directions
of space.

% --- source: tex/s15_shifted#0 ((head))


% --- source: tex/s15_shifted#1 (The shifted Lagrangian in a positive form)
\section{The shifted Lagrangian in a positive form}
\label{sec:shifted}

Everything in the arena so far has been kinematics. The arena supplies a candidate first-order shift operator, built from the conversion
between times and lengths (\cref{sec:clock}), and this part adopts a shifted quadratic action for it and asks what motion it generates. For the
adopted shifted quadratic action, the comparison concerns the choice of the quadratic form, positive definite or the split form $\eta$ of the
arena itself, and the two answers are opposite. This section treats the positive form, where the verdict is negative and
sharp; \cref{sec:neutral} treats the split form, where bound motion appears exactly in neutral signature.

% --- source: tex/s15_shifted#2 (Momentum and force)
\subsection{Momentum and force}

Let $X(\tau)\in\R^N$ be an interior configuration, $g$ a constant real $N\times N$ matrix, and
\begin{equation}
\label{eq:alpha}
\alpha_\pm:=\partial_\tau\pm\kappa g,\qquad \mathcal L=\frac m2\,\bigl|\alpha_+X\bigr|^2=\frac m2\,\bigl|\dot X+\kappa gX\bigr|^2 .
\end{equation}
Only the product $\kappa g$ appears: $\kappa=c/\ell$ carries the conversion of \eqref{eq:kappa} and $g$ its direction. The Lagrangian is the
square of a single first-order operator, which is what makes it the natural object to ask Clifford questions of.

\begin{proposition}[Momentum and force]
\label{prop:force}
The canonical momentum and the Euler--Lagrange equation are
\begin{equation}
\label{eq:pdot}
p=\frac{\partial\mathcal L}{\partial\dot X}=m\,\alpha_+X,\qquad \dot p=\kappa\,g^{\!\top}p ,
\end{equation}
and in second-order form, with $a:=\tfrac12(g-g^{\!\top})$,
\begin{equation}
\label{eq:Eg}
\mathcal E_g:=m\bigl(\ddot X+2\kappa a\dot X-\kappa^2g^{\!\top}gX\bigr)=0 .
\end{equation}
\end{proposition}

\begin{proof}
$\partial\mathcal L/\partial\dot X=m(\dot X+\kappa gX)$ and $\partial\mathcal L/\partial X=\kappa g^{\!\top}p$. Substituting $p=m(\dot X+\kappa gX)$ into
$\dot p=\kappa g^{\!\top}p$ gives $m(\ddot X+\kappa g\dot X)=\kappa mg^{\!\top}\dot X+\kappa^2mg^{\!\top}gX$, which rearranges to \eqref{eq:Eg}.
\end{proof}

The force law is $\dot p=\kappa g^{\!\top}p$. It takes the symmetric-looking form $\dot p=\kappa gp$ exactly when $g$ is symmetric, and the
next two results show that this is not a cosmetic condition.

\begin{theorem}[Factorisation]
\label{thm:factor}
$\mathcal E_g=m\,\alpha_-\alpha_+X$ identically if and only if $a=0$.
\end{theorem}

\begin{proof}
For constant $g$, $\alpha_-\alpha_+X=\ddot X-\kappa^2g^2X$, so $\mathcal E_g-m\alpha_-\alpha_+X=m\bigl(2\kappa a\dot X+\kappa^2(g^2-g^{\!\top}g)X\bigr)$. The two
terms involve $\dot X$ and $X$ independently, and both vanish for all trajectories if and only if $a=0$, which makes $g$ symmetric and
$g^2=g^{\!\top}g$.
\end{proof}

\begin{proposition}[The factorised branch is an inverted oscillator]
\label{prop:inverted}
The Noether energy
\begin{equation}
\label{eq:energy}
E=\langle p,\dot X\rangle-\mathcal L=\frac m2|\dot X|^2-\frac{m\kappa^2}{2}\,|gX|^2
\end{equation}
is conserved for every $g$. When $a=0$ the motion is $m\ddot X=-\nabla V$ with $V=-\tfrac12m\kappa^2|gX|^2\le0$, and the initial data whose
solutions stay bounded for $\tau\ge0$ form a linear subspace of dimension exactly $N$ in the $2N$-dimensional phase space; for all real
$\tau$ only the zero-frequency sector, of dimension $\dim\ker g$, stays bounded.
\end{proposition}

\begin{proof}
Writing $u=\alpha_+X$, $E=\tfrac m2|u|^2-m\kappa\langle u,gX\rangle$, and expanding $u$ gives \eqref{eq:energy}. On shell
$\dot E=\tfrac\kappa m\langle p,(g^{\!\top}-g)p\rangle=-\tfrac{2\kappa}m\langle p,ap\rangle=0$ because $a$ is antisymmetric. With $a=0$ the equation is
$\ddot X=\kappa^2g^2X$, and on the eigenplane of an eigenvalue $\lambda\neq0$ of $g$ the solutions are $Ae^{\kappa|\lambda|\tau}+Be^{-\kappa|\lambda|\tau}$,
bounded for $\tau\ge0$ iff $A=0$; on $\lambda=0$ they are $x_0+v_0\tau$, bounded iff $v_0=0$. In each plane one condition is imposed on two
dimensions.
\end{proof}

The factorisation condition and the inverted-potential condition are the same condition. Since $g^{\!\top}g\succeq0$, the potential is
never a well, and the bounded set has measure zero: no open set of initial data stays bounded. Only the velocity-coupled term $2\kappa
a\dot X$ could bind, and it is exactly the term that obstructs factorisation.

\begin{example}[The conversion operator as the shift]
\label{ex:convshift}
Take $g=B_C$, the mixed bivector of \cref{thm:spectrumBC} acting on $S$, which is symmetric. The equation of motion is $\ddot X=\kappa^2B_C^2X$, and
by \cref{thm:spectrumBC} the growth rates are $\kappa|\eta_1\sigma_1+\eta_2\sigma_2+\eta_3\sigma_3|$ for the four sign patterns $\eta$. For the
conversion along a single clock, the symmetric one or any other, $C$ has singular values $(1,0,0)$, every rate equals $\kappa$, and the dynamics cannot see the Clifford structure
(\cref{ex:blindclock}). For $C$ of rank two with $\sigma=(2,1,0)$ the rates are $3\kappa$ and $\kappa$, each on a four-dimensional subspace; for
$\sigma=(1,1,0)$ they are $2\kappa$ and $0$, and the second is a flat direction. Raising the rank makes the instability anisotropic and can
create flat directions, never an oscillation.
\end{example}

% --- source: tex/s15_shifted#3 (No recurrence at nonzero momentum)
\subsection{No recurrence at nonzero momentum}

It is tempting to stop here and blame the tension between factorisation and binding. That would be too generous to the unfactorised
branch. A single identity settles the question for every $g$.

\begin{theorem}[Virial identity]
\label{thm:virial}
With the positive form of \eqref{eq:alpha}, along every solution and for every constant $g$,
\begin{equation}
\label{eq:virial}
\frac{d}{d\tau}\langle X,p\rangle=\frac{|p|^2}{m}\ \ge\ 0 .
\end{equation}
Consequently: (i) no solution with $p(0)\neq0$ is periodic or recurrent; (ii) purely imaginary characteristic exponents may occur, but never
semisimply, so the motion they carry is secular, and the Jordan blocks need not have size two; (iii) a trajectory bounded in phase space has
$p\in L^2$, so either $p\equiv0$ or $p$ decays exponentially and the orbit is not recurrent.
\end{theorem}

\begin{proof}
$\tfrac{d}{d\tau}\langle X,p\rangle=\langle\dot X,p\rangle+\langle X,\kappa g^{\!\top}p\rangle=\langle\dot X+\kappa gX,p\rangle=|p|^2/m$. A recurrent orbit would
return $\langle X,p\rangle$ to an earlier value, which \eqref{eq:virial} forbids unless $p\equiv0$, and $p=e^{\kappa g^{\!\top}\tau}p(0)$ vanishes
identically only if $p(0)=0$. A semisimple imaginary exponent would supply a periodic solution with $p\neq0$. For (iii), boundedness makes
$\langle X,p\rangle$ bounded and monotone, hence convergent, so $\int_0^\infty|p|^2<\infty$, and a bounded square-integrable solution of a
constant-coefficient linear system is zero or exponentially decaying.
\end{proof}

\begin{example}[A Jordan block of size four]
\label{ex:jordan4}
Let $J_2$ be the standard $2\times2$ rotation generator and $g=\bigl(\begin{smallmatrix}J_2&\id_2\\0&J_2\end{smallmatrix}\bigr)$ on $\R^4$, with
$\kappa=m=1$. The eight-dimensional phase-space generator $M$ has its spectrum on $\pm i$, and
$\rank\bigl((M^2+\id)^k\bigr)=6,4,2,0$ for $k=1,2,3,4$. The blocks therefore have size four and the generic solution grows like $\tau^3$, although
every exponent is purely imaginary.
\end{example}

The zero-momentum sector is not empty, and it does oscillate. The general solution is
\begin{equation}
\label{eq:gensol}
p(\tau)=e^{\kappa g^{\!\top}\tau}p_0,\qquad X(\tau)=e^{-\kappa g\tau}\Bigl[X_0+\frac1m\int_0^\tau e^{\kappa gs}e^{\kappa g^{\!\top}s}p_0\,ds\Bigr],
\end{equation}
and the kernel of $\alpha_+$, $X(\tau)=e^{-\kappa g\tau}X_0$, has $p\equiv0$. The integral collapses to $\tau p_0$ only when $g^{\!\top}=-g$. For
$g=J$ antisymmetric with $J^2=-\id$, the zero-momentum solutions are bounded rotations of period $2\pi/|\kappa|$, and any nonzero momentum
adds the secular term $\tau\,e^{-\kappa J\tau}p_0/m$: the orbit becomes an outward spiral (\cref{fig:interior}, left). In that marginal case
$g^{\!\top}g=\id$, $a=J$, the equation reads $\ddot X+2\kappa J\dot X-\kappa^2X=0$, and in the eigenbasis of $J$ the characteristic polynomial is
$(\lambda+i\kappa)^2$, a double root on the imaginary axis; the corresponding Jordan block is exactly the instability just described.

So with a positive form, oscillation exists only on a measure-zero sector and is destroyed by any momentum. What is missing is neither
curvature nor a dependence of $g$ on $\tau$. It is the positive form itself, and the next section removes it.

% --- source: tex/s15_shifted#4 (Towards a stress--energy tensor)
\subsection{Towards a stress--energy tensor}

The worldline statement lifts to a curved metric in two ways, which together indicate where a stress--energy tensor would come from.
First, replacing $\partial_\tau$ by $D/D\tau$ along a worldline in $(M,g)$ and promoting $g$ to a section $G$ of $\End(TM)$, the variation of
$\int\tfrac m2g_{AB}(\alpha_+X)^A(\alpha_+X)^B\,d\tau$ gives $Dp_B/D\tau=\kappa G^A{}_Bp_A$, and the second-order form acquires one new term,
\[
\frac{D^2X}{D\tau^2}+2\kappa a\frac{DX}{D\tau}+\kappa(D_\tau G)X-\kappa^2G^{\!\top}GX=0 ;
\]
the factorisation computation goes through verbatim, since the extra term appears on both sides, so $a=0$ remains the factorisation
condition. Second, for a field $X$ in a bundle with fibre metric $h$ and $D_M=\nabla_M+\kappa G_M$, the action
$\int\sqrt{|g|}\,\tfrac m2h_{AB}g^{MN}(D_MX)^A(D_NX)^B$ has, when $g$ is varied with the fibre metric $h$ and the bundle connection held fixed
(a connection that depends on $g$, such as the Levi-Civita connection on $TM$, needs further terms),
\[
T_{MN}=m\Bigl[h_{AB}(D_MX)^A(D_NX)^B-\tfrac12g_{MN}\,h_{AB}\,g^{PQ}(D_PX)^A(D_QX)^B\Bigr],
\]
whose trace $m(1-\tfrac D2)|DX|^2$ vanishes identically for arbitrary fields only in $D=2$, the same dimension in which the projector sigma model of \cref{sec:ym} is
renormalisable. This is the stress tensor of the shifted kinetic term only, with $G$ a background, and \cref{thm:virial} is a warning about
it: a matter sector whose only oscillations sit at zero momentum, with an energy unbounded below, is not yet a source one can couple to
gravity.

% --- source: tex/s16_neutral#0 ((head))


% --- source: tex/s16_neutral#1 (The split form and binding in neutral signature)
\section{The split form and binding in neutral signature}
\label{sec:neutral}

The norm in \eqref{eq:alpha} is positive definite, and that is what makes $g^{\!\top}g\succeq0$ in \cref{prop:inverted} and $|p|^2\ge0$ in
\cref{thm:virial}. But the arena does not hand us a positive form; it hands us $\eta$, of signature $(3,3)$. Redoing the derivation with
$\eta$ changes one inequality, and with it the verdict.

% --- source: tex/s16_neutral#2 (The split-form Lagrangian)
\subsection{The split-form Lagrangian}

Write $G^\dagger:=\eta^{-1}G^{\!\top}\eta$ for the $\eta$-adjoint and take
\begin{equation}
\label{eq:splitL}
\mathcal L=\frac m2\,\eta\bigl(\alpha_+X,\alpha_+X\bigr),\qquad p=m\,\eta\,(\alpha_+X),\qquad \pi=\eta^{-1}p,\qquad \dot p=\kappa\,G^{\!\top}p ,
\end{equation}
with second-order form and energy
\begin{equation}
\label{eq:splitEg}
\ddot X+2\kappa a\dot X-\kappa^2G^\dagger GX=0,\qquad a=\tfrac12(G-G^\dagger),\qquad E=\frac m2\eta(\dot X,\dot X)-\frac{m\kappa^2}2\eta(GX,GX).
\end{equation}
Everything in \cref{sec:shifted} survives with the transpose replaced by the $\eta$-adjoint: factorisation holds exactly when $a=0$, and $E$
is conserved because $a$ is $\eta$-antisymmetric. What does not survive is the sign of the virial identity.

\begin{theorem}[The split virial identity]
\label{thm:splitvirial}
Let $X$ be a section along a path in a manifold with metric of signature $(3,3)$ and let $D/D\tau$ come from any metric-compatible connection,
torsion included. Then for every $G$,
\begin{equation}
\label{eq:splitvirial}
\frac{d}{d\tau}\,\eta(X,\pi)=\frac{\eta(\pi,\pi)}m .
\end{equation}
Consequently every closed orbit satisfies $\oint\eta(\pi,\pi)\,d\tau=0$: a bound orbit has null momentum on average.
\end{theorem}

\begin{proof}
Note that $\eta(X,\pi)=\langle p,X\rangle$. Metric compatibility gives $\tfrac{d}{d\tau}\eta(X,\pi)=\eta(DX/D\tau,\pi)+\eta(X,D\pi/D\tau)$, and torsion does
not enter this Leibniz rule. Substituting $DX/D\tau=\pi/m-\kappa GX$ and $D\pi/D\tau=\kappa G^\dagger\pi$ gives $\eta(\pi,\pi)/m-\kappa\,\eta(GX,\pi)+\kappa\,\eta(X,G^\dagger\pi)$,
and the last two terms cancel by the definition of the adjoint.
\end{proof}

Neither curvature nor torsion appears in \eqref{eq:splitvirial}; the identity sees only $\eta$. In definite signature $\eta(\pi,\pi)$ has
one sign and the identity forbids recurrence, which is \cref{thm:virial}. In split signature it does not, and the question becomes whether
solutions with null momentum on average exist. They do, and the operator that produces them is one we already have.

% --- source: tex/s16_neutral#3 (Binding exactly in neutral signature)
\subsection{Binding exactly in neutral signature}

\begin{theorem}[Factorisation and binding coexist exactly in neutral signature]
\label{thm:neutral}
Let $\eta$ have signature $(P,Q)$ on $\R^N$. There is a linear map $G$ with
\begin{enumerate}[label=(\roman*),leftmargin=2.2em]
\item $G=G^\dagger$, so that the factorisation of \cref{thm:factor} holds, and
\item $G^2=-\id$, so that \eqref{eq:splitEg} reads $\ddot X=-\kappa^2X$ and every solution is periodic of period $2\pi/\kappa$,
\end{enumerate}
if and only if $P=Q$. Such a $G$ is an anti-isometry, $\eta(GX,GY)=-\eta(X,Y)$, and on $\R^{3,3}$ one of them, maximal in that it uses every
paired direction, is the binding root of \cref{prop:twoK}:
\begin{equation}
\label{eq:Gpair}
G:\ e_i\longmapsto e_{i+3},\qquad e_{i+3}\longmapsto-e_i\qquad(i=1,2,3).
\end{equation}
\end{theorem}

\begin{proof}
Given (i) and (ii), $\eta(GX,GY)=\eta(X,G^\dagger GY)=\eta(X,G^2Y)=-\eta(X,Y)$, so $G$ carries a maximal positive subspace isometrically onto a
negative one and back, which forces $P=Q$. Conversely \eqref{eq:Gpair} satisfies $G^2=-\id$, and
$\eta(Ge_i,e_{i+3})=\eta(e_{i+3},e_{i+3})=-1=\eta(e_i,-e_i)=\eta(e_i,Ge_{i+3})$, so $G=G^\dagger$. With $G^\dagger G=G^2=-\id$, \eqref{eq:splitEg}
has $a=0$ and becomes $\ddot X=-\kappa^2X$.
\end{proof}

The linear algebra is inherited: a complex structure that is anti-isometric for $\eta$ is a Norden, or anti-K\"ahler, structure, and its
existence forces neutral signature \cite{BorowiecFrancavigliaVolovich00}. What the arena adds is the reading of \eqref{eq:Gpair} as the binding
root of the pairing.

This reverses the verdict of \cref{sec:shifted} at its root. In definite signature the factorised branch is an inverted oscillator and
nothing binds; in neutral signature the same factorisation condition $G=G^\dagger$ admits operators of purely imaginary spectrum, and the
maximal one makes the whole interior a harmonic well of one frequency. The obstruction of \cref{thm:virial} was a theorem about the
positive form, not about the arena. Two qualifications keep the statement honest. The normalisation $G^2=-\id$ imposes one common
frequency; any $\eta$-self-adjoint shift built from independent paired blocks $\mu_iJ_2$ with distinct $\mu_i>0$ also factorises and binds, at
unequal frequencies, so the construction is a family and $G^2=-\id$ picks a member. And no equation of motion for $G$ has been written:
$G$ is adopted, and calling its consequences a selection would be premature.

\begin{proposition}[How much can bind]
\label{prop:howmuch}
Let $G=G^\dagger$ be diagonalisable over $\C$ and $\kappa\neq0$. An eigenvector $v$ with non-real eigenvalue satisfies $\eta(v,\bar v)=0$, and
eigenvectors of different eigenvalues are $\eta$-orthogonal. The real span of the eigenvectors with non-real eigenvalues is therefore the
$\eta$-orthogonal complement of the real eigenspaces; it is nondegenerate of signature $(k',k')$. Let $V_{\rm osc}$ be its part belonging to the
nonzero purely imaginary eigenvalues. $V_{\rm osc}$ is nondegenerate of signature $(k,k)$ with $k\le\min(P,Q)$, $\eta$-orthogonal to the remaining
spectral subspaces, and every solution of $\ddot X=\kappa^2G^2X$ with position and velocity in $V_{\rm osc}$ is bounded (quasi-periodic; periodic when
the frequencies are commensurate). Zero eigenvalues give constant solutions and linear drift, and eigenvalues with nonzero real part give no
nonzero solution bounded for all times. On Minkowski space $\R^{1,3}$ at most a single two-plane can bind; on $\R^{3,3}$ all six directions can,
simultaneously and at one frequency.
\end{proposition}

A non-real spectrum is not enough. For $\eta=\operatorname{diag}(1,-1)$ the operator $G=\bigl(\begin{smallmatrix}1&-1\\1&1\end{smallmatrix}\bigr)$ is
$\eta$-self-adjoint and diagonalisable with eigenvalues $1\pm i$, and $X(\tau)=e^{\kappa\tau}(\cos\kappa\tau,\sin\kappa\tau)$ solves $\ddot X=\kappa^2G^2X$:
the oscillation is exponentially modulated, and no nonzero solution of this block is bounded on $\R$. Hence the restriction to the purely
imaginary spectrum, $\operatorname{spec}G\subset\R\cup i\R$ being the case in which the two subspaces agree (a3-11, exact).

The statement concerns the whole subspace $V_{\rm osc}$. Let $\eta=\bigl(\begin{smallmatrix}0&K\\K&0\end{smallmatrix}\bigr)$ on $\R^4$, with
$K=\operatorname{diag}(1,-1)$. For the diagonalisable $G=\operatorname{diag}(J_2,J_2)$, which is $\eta$-self-adjoint, the plane of one conjugate
pair is totally null, although the whole of $V_{\rm osc}$ has signature $(2,2)$. For the non-diagonalisable
$G=\bigl(\begin{smallmatrix}J_2&\id\\0&J_2\end{smallmatrix}\bigr)$ even the subspace statement fails: $G$ is $\eta$-self-adjoint and its bounded
motions lie in a plane on which $\eta$ vanishes. The conclusions used later concern $G^2=-\id$, which is
diagonalisable, and stand.

\begin{example}[One pairing plane]
\label{ex:oneplane}
In the plane $(e_1,e_4)$ write $X=x\,e_1+y\,e_4$, so that $G$ acts as $\bigl(\begin{smallmatrix}0&-1\\1&0\end{smallmatrix}\bigr)$ and $z=x+iy$. The
general solution of $\ddot X=-\kappa^2X$ is $z(\tau)=c_-e^{-i\kappa\tau}+c_+e^{i\kappa\tau}$, an ellipse with semi-axes $|c_-|+|c_+|$ and
$\bigl||c_-|-|c_+|\bigr|$. The momentum $\alpha_+X=\dot X+\kappa GX$ corresponds to $\dot z+i\kappa z=2i\kappa\,c_+e^{i\kappa\tau}$, so it is carried
entirely by the counter-rotating part, and the zero-momentum orbits are the circles $z=c_-e^{-i\kappa\tau}$, traversed at angular velocity
$-\kappa$. The virial quantity $\eta(\pi,\pi)=m^2\eta(\alpha_+X,\alpha_+X)=m^2\Re\bigl((2i\kappa c_+e^{i\kappa\tau})^2\bigr)$ oscillates with mean
zero, as \cref{thm:splitvirial} requires of a closed orbit.
\end{example}

\begin{figure}[t]
\centering
\begin{tikzpicture}
\begin{axis}[name=left,width=6.8cm,height=6.8cm,axis equal image,xmin=-5.4,xmax=5.4,ymin=-5.4,ymax=5.4,
  axis lines=middle,xtick=\empty,ytick=\empty,xlabel={$x$},ylabel={$y$},
  title={\small positive form, $g=J$}]
\addplot[gray,dashed,thick,domain=0:360,samples=90] ({cos(x)},{sin(x)});
\addplot[scol,thick,domain=0:5*pi,samples=600,variable=\t] ({cos(deg(\t))+0.3*\t*sin(deg(\t))},{-sin(deg(\t))+0.3*\t*cos(deg(\t))});
\node[scol,anchor=north west,font=\small] at (axis cs:-5.3,5.3) {$p_0\neq0$: spiral};
\node[gray,anchor=south west,font=\small] at (axis cs:-5.3,-5.3) {dashed: $p=0$};
\end{axis}
\begin{axis}[at={(left.east)},anchor=west,xshift=1.0cm,width=6.8cm,height=6.8cm,axis equal image,xmin=-2.0,xmax=2.0,ymin=-2.0,ymax=2.0,
  axis lines=middle,xtick=\empty,ytick=\empty,xlabel={$x$},ylabel={$y$},
  title={\small split form, $G$ the pairing}]
\addplot[gray,dashed,thick,domain=0:360,samples=90] ({cos(x)},{sin(x)});
\addplot[tcol,thick,domain=0:360,samples=180] ({1.4*cos(x)},{-0.6*sin(x)});
\addplot[gcol,very thick,domain=0:360,samples=180] ({0.85*cos(x)-0.85*sin(x)},{0.85*cos(x)-0.85*sin(x)});
\node[tcol,anchor=north east,font=\small] at (axis cs:1.95,-0.7) {$|c_+|<|c_-|$};
\node[gcol,anchor=north west,font=\small] at (axis cs:-1.95,-1.45) {$|c_+|=|c_-|$};
\node[gray,anchor=south east,font=\small] at (axis cs:-0.55,0.95) {$p=0$};
\end{axis}
\end{tikzpicture}
\caption{Interior motion in one pairing plane (\cref{ex:oneplane}). Left: with the positive form and an antisymmetric shift the
zero-momentum orbits are circles, and any momentum turns them into outward spirals, as \cref{thm:virial} requires. Right: with the split form
and the binding root $G$ every orbit is closed; the zero-momentum orbits are circles, the momentum-carrying ones ellipses, degenerating to a
segment when the two rotating parts have equal size.}
\label{fig:interior}
\end{figure}

% --- source: tex/s16_neutral#4 (The orbits as duality rotations)
\subsection{The orbits as duality rotations}

Read through the field-strength dictionary of \cref{thm:emfield}, an interior configuration $X$ is a two-form and the binding root is the
Hodge star. The bound motion then has a transparent meaning.

\begin{proposition}[Bound orbits in field-strength language]
\label{prop:orbitsEM}
Let $G$ be the pairing \eqref{eq:Gpair} and $z=\bB+i\bE$.
\begin{enumerate}[label=(\roman*),leftmargin=2.2em]
\item The zero-momentum orbits are $X(\tau)=e^{-\kappa G\tau}X_0$, that is $z(\tau)=e^{-i\kappa\tau}z_0$: continuous duality rotations, along which
$q(z(\tau))=e^{-2i\kappa\tau}q(z_0)$. A non-null orbit is a single spacetime plane at exactly four instants per period, alternately of magnetic
and electric type, and a null orbit is a single plane at every instant.
\item Every solution is $z(\tau)=c_-e^{-i\kappa\tau}+c_+e^{i\kappa\tau}$ with $c_\pm\in\C^3$, $\alpha_+X$ corresponds to $2i\kappa c_+e^{i\kappa\tau}$ (the momentum field $\pi=m\alpha_+X$ to $m$ times it), and the
period averages are
\[
\langle q\rangle=2\,c_-\cdot c_+,\qquad \langle\bE\cdot\bB\rangle=\Im(c_-\cdot c_+),\qquad \langle|\bB|^2-|\bE|^2\rangle=2\Re(c_-\cdot c_+).
\]
\item The energy is $\mathcal E=2m\kappa^2\Re(c_-\cdot c_+)$, and the complex quantity $\tfrac m2\bigl(q(\dot z)+\kappa^2q(z)\bigr)=2m\kappa^2\,c_-\cdot c_+$ is
conserved; its imaginary part is a second constant of motion.
\item On the orbits with $c_+=\bar c_-$ the vector $z$ is real at every instant: $\bE=0$ and $\bE\cdot\bB=0$ throughout, while the momentum is not
zero.
\end{enumerate}
\end{proposition}

On a zero-momentum orbit $\bE\cdot\bB(\tau)=\tfrac12|q_0|\sin(\phi_0-2\kappa\tau)$ with $q_0=|q_0|e^{i\phi_0}$, and the rest-frame field of
\cref{sec:orth} turns from purely magnetic through screws to purely electric and back, at the fixed rest-frame energy density
$\tfrac12|q_0|$. So the binding flow is what breaks orthogonality on every non-null orbit without momentum, while orbits that carry momentum
can stay orthogonal throughout. The mean non-orthogonality and the energy of a general orbit are the imaginary and real parts of one conserved
complex number $c_-\cdot c_+$, and neither determines the other.

\begin{theorem}[The parity-odd term is a total derivative]
\label{thm:EBsplit}
For the split-form Lagrangian with the pairing,
\[
\mathcal L_\kappa=\frac m2\Bigl(\eta(\dot X,\dot X)-\kappa^2\eta(X,X)\Bigr)-m\kappa\,\frac{d}{d\tau}\bigl(\bE\cdot\bB\bigr).
\]
Consequently the cross term has identically vanishing Euler--Lagrange expression, the equation $\ddot X=-\kappa^2X$ comes from the first bracket
alone, and the momenta of $\mathcal L_\kappa$ and $\mathcal L_{-\kappa}$ differ by $2m\kappa\,\eta GX$. Parity maps $\mathcal L_\kappa$ to $\mathcal L_{-\kappa}$
and exchanges the orbits without momentum with those carried entirely by the counter-rotating part, while $C$ and $T$ leave $\mathcal L_\kappa$
invariant.
\end{theorem}

\begin{proof}
$G$ is $\eta$-self-adjoint, so $\tfrac{d}{d\tau}\eta(X,GX)=2\eta(\dot X,GX)$, and $\eta(X,GX)=-2\,\bE\cdot\bB$ by \cref{thm:emfield}; $\eta(GX,GX)=-\eta(X,X)$
because $G$ is an anti-isometry. Expanding $\eta(\dot X+\kappa GX,\dot X+\kappa GX)$ gives the formula. Parity reverses $\bE$ and fixes $\bB$, so
$PGP^{-1}=-G$ and $P$ acts on $z$ as complex conjugation.
\end{proof}

The binding term carries the one parity-odd ingredient of the interior, the sign of $\kappa$, and it sits entirely in a total derivative.
Since $\bE\cdot\bB$ is a single-valued function on the configuration space, the canonical transformation $p\mapsto p-m\kappa\,\eta GX$ removes it,
classically and as a phase of wave functions; the shift by $2m\kappa\,\eta GX$ of \cref{thm:EBsplit} relates the two gauges $\pm\kappa$. The structure is that of the $\theta$-term of electrodynamics, a parity-odd coefficient times a
total derivative, with the difference that here it has no topological sector in which to act. In the positive form the same cross term is
gyroscopic, with Euler--Lagrange expression $2m\kappa G\dot X$, and the sign of $\kappa$ enters the equations of motion; that is another way of
seeing why the two forms behave so differently.

\subsection{Frame field and momentum field}
\label{sec:framemomentum}

Two fields can be read along an interior solution: the frame field $X$ itself and the momentum field $\pi=\eta^{-1}p=m\,\alpha_+X$.
For a field point $Y$ with $z(Y)=\bB+i\bE$, $E_6(Y)$, $E_0(Y)$ and $\mathbf S(Y)$ are its undivided energy, its rest energy and its
Poynting vector, and $h(z,w)=\sum_j\bar z_jw_j$ (\cref{en:defs}). By \cref{prop:orbitsEM}, every solution of $\ddot X=-\kappa^2X$ is $z(\tau)=c_-e^{-i\kappa\tau}+c_+e^{i\kappa\tau}$,
$c_\pm\in\C^3$, the momentum field $\pi=m\alpha_+X$ corresponds to $2im\kappa c_+e^{i\kappa\tau}$, the Noether energy is $\mathcal E=2m\kappa^2\Re(c_-\cdot c_+)$, and the
complex quantity $\mathcal E_\C:=\tfrac m2\bigl(q(\dot z)+\kappa^2q(z)\bigr)=2m\kappa^2\,c_-\cdot c_+$ is conserved. The pointwise energies of the frame
field and of the momentum field behave differently along these solutions.

\begin{proposition}[Frame field and momentum field in the interior]
\label{en:interior}
Along every solution of the split-form interior with the binding root:
\begin{enumerate}[label=(\roman*),leftmargin=2.2em]
\item The momentum field is $z(\pi)=2im\kappa\,c_+e^{i\kappa\tau}$, that is $\pi(\tau)=e^{\kappa\tau G}\pi_0$: it moves by pure duality, and
$E_6(\pi)=\tfrac12|p|^2=2m^2\kappa^2|c_+|^2$, $E_0(\pi)=2m^2\kappa^2|c_+\cdot c_+|$ and $\mathbf S(\pi)$ are constant.
\item The frame field has
\[
|z|^2=|c_-|^2+|c_+|^2+2\Re\bigl(h(c_-,c_+)e^{2i\kappa\tau}\bigr),\qquad q(z)=e^{-2i\kappa\tau}c_-\cdot c_-+2c_-\cdot c_++e^{2i\kappa\tau}c_+\cdot c_+ ,
\]
so its $E_6$ is constant iff $h(c_-,c_+)=0$, and its $E_0$ oscillates in general: $c_-=e_1$, $c_+=e_2$ give $\mathcal E=0$, $E_6=1$ and
$E_0=|\cos2\kappa\tau|$.
\item Both pure senses, $z=e^{-i\kappa\tau}c_-$ and $z=e^{i\kappa\tau}c_+$, are duality rotations of the frame field, with $\mathcal E_\C=0$ and with
constant $E_6$ and $E_0$ of the frame field. On the first, the orbits without momentum, $p=0$ while $E_6=\tfrac12|c_-|^2>0$; on the second
$|p|^2=8m^2\kappa^2E_6$.
\item The parameter $m$ is an overall factor of the Lagrangian and does not enter $\ddot X=-\kappa^2X$.
\end{enumerate}
\end{proposition}

\begin{proof}
$\dot z+i\kappa z=2i\kappa c_+e^{i\kappa\tau}$ is $z(\alpha_+X)$, since $G$ acts as $i$, and $\pi=m\alpha_+X$; with $|\pi|=|p|$ ($\eta$ is orthogonal)
this gives (i), and $\overline{z(\pi)}\times z(\pi)$, $|z(\pi)|$ and $|z(\pi)\cdot z(\pi)|$ are invariant under multiplication of $z(\pi)$ by a phase. (ii) is the
expansion of $\bar z\cdot z$ and $z\cdot z$; in the example $|z|^2=2$ and $q=2\cos2\kappa\tau$. (iii) is (i) and (ii) at $c_+=0$ or $c_-=0$, with
$|p|^2=|z(\pi)|^2=4m^2\kappa^2|c_+|^2$. (iv) is read off \eqref{eq:splitL} and \cref{thm:neutral}.
\end{proof}

Item (i) and the first half of (iii) make explicit what \cref{prop:orbitsEM} states about the momentum and about the orbits without momentum;
item (ii) and the second half of (iii) are fixed by the algebra and complete the picture, in which parity exchanges the two senses
(\cref{thm:EBsplit}). The
quantities the interior conserves belong to the momentum field: $E_6(\pi)$, $E_0(\pi)$ and $\mathbf S(\pi)$ are constant already in the clamped
interior, because $\pi$ moves by a duality rotation. The same functions of the frame field are not conserved, and $E_6$ of the frame field is not a
function of $|p|^2$: on the orbits without momentum $|p|^2=0$ while $E_6>0$, on the duality rotations of the other sense $|p|^2=8m^2\kappa^2E_6$.
Duality acts along both pure senses, so duality orbits are not the orbits without momentum; that sense is singled out by the momentum, not by
duality. And the interior's conserved energy is not undivided: it is the real part of a constant of duality charge two, carried by neither sense
alone.

% --- source: tex/s16_neutral#5 (What the binding root selects, and what tears it)
\subsection{What the binding root selects, and what tears it}

The binding root is not merely available; it carries structure. By \cref{thm:GLorentz} its centraliser is a Lorentz algebra $\lG$, not
conjugate to the interior one (\cref{prop:notconj}). The same operator limits what transport can do to the bound sector.

\begin{proposition}[Transport and a homogeneous threshold]
\label{prop:tear}
Let a connection along the path take values $\omega$ in $\so{3,3}$. If $\omega$ lies in the compact part of $\lG$, transport and the shift commute
and the orbit is carried rigidly. If $\omega$ has a boost component of strength $c$ in a pairing plane $(e_a,e_{a+3})$, the homogeneous $2\times2$
block there has exponents $\pm\sqrt{c^2-\kappa^2}$: oscillatory for $c<\kappa$ and real for $c>\kappa$.
\end{proposition}

\begin{proof}
In the basis $(e_a,e_{a+3})$ the boost generator is $\sigma=\bigl(\begin{smallmatrix}0&1\\1&0\end{smallmatrix}\bigr)$ and $G$ is
$\bigl(\begin{smallmatrix}0&-1\\1&0\end{smallmatrix}\bigr)$, so $-c\sigma+\kappa G=\bigl(\begin{smallmatrix}0&-c-\kappa\\-c+\kappa&0\end{smallmatrix}\bigr)$, whose
square is $(c^2-\kappa^2)\id$.
\end{proof}

The threshold classifies exponents of the homogeneous block, not boundedness of the full flow, and the difference is real. With the momentum
included, the phase-space generator $M=\bigl(\begin{smallmatrix}-c\sigma-\kappa G&\id/m\\0&-c\sigma+\kappa G\end{smallmatrix}\bigr)$ satisfies
$M^2-(c^2-\kappa^2)\id=\bigl(\begin{smallmatrix}0&-2c\sigma/m\\0&0\end{smallmatrix}\bigr)$, nilpotent and nonzero for $c\neq0$. At $\kappa=m=1$ and
$c=\tfrac12$ one finds $\det(\lambda\id-M)=(\lambda^2+\tfrac34)^2$ with $\rank(M^2+\tfrac34\id)=2$: Jordan blocks, and generic secular growth strictly
below the threshold. Along an open path, moreover, a connection coefficient can be changed by a frame rotation, so growth on one worldline
proves nothing; the invariant is the holonomy around a closed path. \Cref{sec:curved} computes it in an explicit curved model and finds that
normal curvature converts zero-momentum orbits into momentum-carrying ones, through exactly the failure of the normal rotation to commute with
$G$.

What this part delivers is a bound sector, a frequency, a centralising Lorentz algebra and an exponent threshold, all from an adopted binding
operator. It does not deliver a spectrum: $\kappa$ is not fixed, the single frequency is degenerate across the three pairs, and the energy is
indefinite, so these are bound orbits and not bound states in any quantum sense. Nor does it settle the relation between the two square roots of
the pairing beyond \cref{lem:Gclifford}: $\cK$ is $\eta$-antisymmetric and grades, $G$ is $\eta$-self-adjoint and binds, and they differ by the
Cartan involution $\Ad_{e_1e_2e_3}$.

\subsection{Bound orbits are not bound states}
\label{sec:noground}

The orbits of \cref{thm:neutral} are classical. Quantised in the most direct way, the same oscillators have no state of lowest energy.

\begin{corollary}[No ground state]
\label{mb:noground}
Quantise the interior oscillators of the full frame canonically, mode by mode, with positive norm. Then the spectrum of the energy
$\mathcal E=\tfrac m2\bigl(\eta(\dot X,\dot X)+\kappa^2\eta(X,X)\bigr)$ is $\hbar\kappa\Z$, every level is infinitely degenerate, and every state has one
of lower energy. Making the spatial energies positive instead requires negative norm.
\end{corollary}

\begin{proof}
Since $G$ is an anti-isometry, $\mathcal E$ is the energy of \eqref{eq:splitEg}, and
$\mathcal E=\sum_a\eta_{aa}\,\tfrac m2(\dot x_a^2+\kappa^2x_a^2)$. The cross term of $\cL_\kappa$ is a total derivative (\cref{thm:EBsplit}); it
shifts the momenta by a canonical transformation and leaves the spectrum unchanged. A spatial mode $x$ has Lagrangian
$-\tfrac m2(\dot x^2-\kappa^2x^2)$, momentum $p_x=-m\dot x$ and Hamiltonian $-(p_x^2/2m+m\kappa^2x^2/2)$. With
$x=\sqrt{\hbar/2m\kappa}\,(\hat a+\hat a^\dagger)$ and $p_x=i\sqrt{\hbar m\kappa/2}\,(\hat a^\dagger-\hat a)$,
$p_x^2/2m+m\kappa^2x^2/2=\tfrac{\hbar\kappa}2(\hat a\hat a^\dagger+\hat a^\dagger\hat a)$, so a temporal mode contributes $\hbar\kappa(j+\half)$ and a spatial one
$-\hbar\kappa(j+\half)$, $j=0,1,\dots$. On $\R^{3,3}$ the half-integers cancel and a state of occupations $(j_1,j_2,j_3;j_4,j_5,j_6)$ has energy
$\hbar\kappa(j_1+j_2+j_3-j_4-j_5-j_6)$. Every integer occurs; adding one temporal and one spatial quantum keeps the energy, and adding a
spatial quantum lowers it by $\hbar\kappa$. A positive spatial energy needs $[\hat a,\hat a^\dagger]=-1$ on the spatial modes, a negative norm.
\end{proof}

The per-mode quantisation is the inherited indefinite oscillator \cite{PaisUhlenbeck50}, and the corollary is an obstruction under the
hypothesis of positive norm. The half-integers cancel because the signature is neutral, which is the condition of \cref{thm:neutral}
under which factorisation and binding coexist: the signature that closes every classical orbit leaves the quantised oscillators without a
ground state. The zero-point values on coordinate subspaces are those of \cref{thm:subfs}.

% --- source: tex/s17_closure#0 ((head))


% --- source: tex/s17_closure#1 (Closing the force law)
\section{Closing the force law}
\label{sec:closure}

The force law $\dot p=\kappa G^{\!\top}p$ carries its rate $\kappa$ as a constant put in by hand. A force law of this shape is familiar from
rigid bodies and ideal fluids, where the equation of motion is a geodesic equation on a Lie group and the rate is itself dynamical. This
section asks whether the interior law can be closed in that sense: whether there is a Lie--Poisson system whose bracket and inertia are built
from the same data that build $G$, and whose solutions contain the force law with $\kappa$ as a dynamical variable. A bracket or an inertia is
called \emph{derived from the Clifford structure} when it is equivariant under the algebra $\lG$ of \cref{thm:GLorentz} and built from
$\eta$, $G$ and the bivectors. The answers, in brief: no closure lives on the frame alone; the closure on one extra direction is unique; its
inertia is constrained by two requirements that exclude each other; and the dynamics it generates is a pendulum in the duality angle that
drifts the frame except on null data.

Throughout, $V=\R^{3,3}$, a covector $p$ is a column with $\langle p,X\rangle=p^{\!\top}X$, $\pi=\eta^{-1}p$, and $z(X)=x+iy$ for
$X=\sum_j(x_je_j+y_je_{j+3})$, so that $G$ acts as $i$. For a covector put
\begin{equation}
\label{eq:wdef}
C(p)=p^{\!\top}\eta^{-1}p,\qquad D(p)=p^{\!\top}G\eta^{-1}p,\qquad w(p):=q\bigl(z(\pi)\bigr)=C(p)-iD(p).
\end{equation}
In the field-strength reading, with $(\bB;\bE)=\pi$, $w=|\bB|^2-|\bE|^2+2i\,\bB\cdot\bE$ is the complex invariant of the field carried by the
momentum.

% --- source: tex/s17_closure#2 (The force law as a dual transport)
\subsection{The force law as a dual transport}

Euler--Arnold equations are geodesic equations on a Lie group with a one-sided invariant metric \cite{Arnold66,ArnoldKhesin98}; for a
left-invariant metric they read $\dot\mu=\ad_w^{\!\top}\mu$ with $\mu=Aw$ for an inertia $A$. The rigid body is the case of $SO(3)$, where this is
$\dot{\mathbf M}=\mathbf M\times\boldsymbol\Omega$ (\cref{sec:rb:euler}), and incompressible Euler flow the case of the volume-preserving diffeomorphisms.

\begin{proposition}[The dictionary]
\label{prop:EAdict}
\begin{enumerate}[label=(\roman*),leftmargin=2.2em]
\item Since $\langle G^{\!\top}p,v\rangle=\langle p,Gv\rangle$, the law $\dot p=\kappa G^{\!\top}p$ transports the covector $p$ by the transpose of $\kappa G$.
It is the Lie--Poisson equation $\dot p=\ad_w^{\!\top}p$ of a Lie bracket on the frame exactly when $\kappa G=\ad_w$ is an inner derivation.
\item For a bracket with trivial centre on the $6n$-dimensional frame the inner derivations fill a $6n$-dimensional subspace of $\so{3n,3n}$, of
dimension $3n(6n-1)$, when $\eta$ is invariant: a fraction $2/(6n-1)$.
\item When $\kappa G=\ad_u$ and $p=Au$, the shift operators $\alpha_\pm=\partial_\tau\pm\kappa G$ are the constrained variations
$\delta u=\dot\xi\pm[u,\xi]$ of Euler--Poincar\'e reduction \cite{MarsdenRatiu99,HolmMarsdenRatiu98}, in the body and in the spatial frame.
\end{enumerate}
\end{proposition}

\begin{proposition}[The binding root is inner for no bracket]
\label{prop:EAnotinner}
For every Lie bracket and every $w$, $\ad_w(w)=0$, so every inner derivation is singular. The binding root satisfies $G^2=-\id$ and is
invertible, so $\kappa G$ is inner for no Lie bracket on $\R^{3n,3n}$, invariant or not, and the force law with that $G$ is never a
Lie--Poisson equation on the frame itself. It is one on the extension $\R\hat g\ltimes\R^{3n,3n}$, the frame taken abelian and $\hat g$ acting on
it by $G$, where it is the Lie--Poisson equation of the fixed element $\kappa\hat g$ restricted to the ideal.
\end{proposition}

\begin{theorem}[Transport and factorisation are complementary]
\label{thm:EAdichotomy}
$\End(\R^{3n,3n})=\so{3n,3n}\oplus\{G:G=G^\dagger\}$, of dimensions $3n(6n-1)+3n(6n+1)=36n^2$. Along $\dot p=\kappa G^{\!\top}p$ the Casimir
$C=\eta^{-1}(p,p)$ changes at the rate $dC/d\tau=2\kappa\,p^{\!\top}G\eta^{-1}p$, which vanishes for every $p$ exactly when $G$ lies in the first
summand, while factorisation in the split form requires the second. The binding root lies in the second summand and moves the Casimir at the rate
$dC/d\tau=-4\kappa\,\bB\cdot\bE$ in the variables $(\bB;\bE)=\eta p$.
\end{theorem}

At the first rung the frame itself is a Lie algebra. With $L_i=e_i$ and $K_i=e_{i+3}$ obeying the brackets of \cref{thm:emfield}, $\R^{3,3}$
becomes $\so{1,3}$ with Killing form $-4\eta$; in the variable $z=\bB+i\bE$ the bracket is the complex cross product, $G$ is multiplication by $i$
and lies in the centroid, $[Gx,y]=G[x,y]$, and $z\cdot z$ is the form $q$. The inner derivations fill six of the fifteen dimensions of
$\so{3,3}$. Closing the force law into an Euler--Arnold equation on the frame would need a bracket for which the shift is inner, and by
\cref{prop:EAnotinner} there is none. A closure must leave the frame.

% --- source: tex/s17_closure#3 (No closure on the frame; the rotor is unique)
\subsection{No closure on the frame; the rotor is unique}

\begin{proposition}[Equivariant maps and forms]
\label{prop:lGhoms}
For $\lG=\operatorname{span}\{L_i,K_i\}$: $\Hom_{\lG}(V,V)=\operatorname{span}\{\id,G\}$ and $\Hom_{\lG}(V,\R)=0$;
$\Hom_{\lG}(\Lambda^2V,V)=\operatorname{span}\{\beta_0,G\beta_0\}$, where $\beta_0(X,Y)$ is the complex cross product $z(X)\times z(Y)$ returned to
$V$, and $\Hom_{\lG}(\Lambda^2V,\R)=0$; and the $\lG$-invariant symmetric forms on $V^*$ are $\operatorname{span}\{\eta^{-1},G\eta^{-1}\}$.
\end{proposition}

\begin{theorem}[The brackets on the frame]
\label{thm:brackets}
The $\lG$-equivariant Lie brackets on $V$ are $\beta_L=L\beta_0$ with $L=a+bG$. The bracket $\beta_0$ makes $(V,\beta_0)\cong\lG$, its derivations
are the six inner ones, its centroid is $\operatorname{span}\{\id,G\}\cong\C$, so $G$ is not a derivation, and $\beta_L$ has Killing form
$-4\bigl[(a^2-b^2)\eta+2ab\,\eta G\bigr]$.
\end{theorem}

\begin{proof}
By \cref{prop:lGhoms} every equivariant antisymmetric bilinear map is $(c_1+c_2G)\beta_0$. Since $L$ lies in the centroid,
$\beta_L(\beta_L(X,Y),Z)=L^2\beta_0(\beta_0(X,Y),Z)$, and $\beta_L$ inherits the Jacobi identity.
\end{proof}

\begin{theorem}[No closure on the frame alone]
\label{thm:noframe}
Let $L=a+bG\neq0$. For every symmetric $N$, with $\mathfrak g^*$ identified with $\mathfrak g$ through $\eta$, the Euler--Arnold field
$\beta_0(M,NM)$ conserves $\eta(M,M)$ and $\eta(M,GM)$, whereas along the force law, in the form $\dot M=\kappa GM$,
$\tfrac{d}{d\tau}\eta(M,M)=2\kappa\,\eta(M,GM)$ and $\tfrac{d}{d\tau}\eta(M,GM)=-2\kappa\,\eta(M,M)$. No Hamiltonian flow on $(V,\beta_L)^*$ coincides
with a nontrivial force law ($\kappa\neq0$) on an open set $U$, and this remains true with a continuous state-dependent rate that does not vanish
identically on $U$.
\end{theorem}

\begin{proof}
The two quadratic forms are the Casimirs of $(V,\beta_L)^*$ and are conserved by every Hamiltonian flow, while the force law rotates them into
each other. Equality on an open set would make both vanish there, that is $z\cdot z=0$ on an open set, which no nonzero polynomial does. For a
state-dependent rate $f$ not identically zero on $U$, continuity gives an open $W\subset U$ on which $f\neq0$, and the same argument applies on $W$.
(On the null cone, where both Casimirs vanish, the equations $\beta_0(M_s,NM_s)=f_s\,GM_s$ at twelve rational null vectors force every $f_s=0$ and
$N\in\operatorname{span}\{\id,G\}$, for which $\beta_0(M,NM)\equiv0$; this sampled computation is a check, not part of the proof.)
\end{proof}

A closure therefore needs a direction beyond the frame. Adjoin one generator $\hat g$, the \emph{rotor}, on which $\lG$ acts trivially; the rate
$\kappa$ is a Lorentz scalar, since $G$ commutes with $\lG$.

\begin{theorem}[Uniqueness of the closure]
\label{thm:closureunique}
Let $\mathfrak s=\R\hat g\oplus V$ carry an $\lG$-equivariant Lie bracket for which $\ad_{\hat g}$ does not vanish on $V$. Then
$[\hat g,X]=(a+bG)X$ and $[X,Y]=0$ for $X,Y\in V$, with $(a,b)\neq(0,0)$. If moreover the rotor rotates the frame covector without dilating it,
as the force law does, then $a=0$ and, after rescaling $\hat g$, $\mathfrak s\cong\R\ltimes_GV$: $[\hat g,X]=GX$, $[X,Y]=0$.
\end{theorem}

\begin{proof}
Equivariance gives $[\hat g,X]=AX+\alpha(X)\hat g$ with $A\in\operatorname{span}\{\id,G\}$ and $\alpha=0$, and $[X,Y]=(c_1+c_2G)\beta_0(X,Y)$, the
$\hat g$-component vanishing by \cref{prop:lGhoms}. The Jacobi identity for $(\hat g,X,Y)$ reads $(a+bG)(c_1+c_2G)\beta_0(X,Y)=0$, a linear system in
$(c_1,c_2)$ with determinant $a^2+b^2\neq0$, so $c_1=c_2=0$. On $\mathfrak s^*$ the rotor acts on $p$ through $(a+bG)^{\!\top}$, and
$\tfrac{d}{d\tau}|p|^2=2a\,(h/I_0)|p|^2$ because $G$ is antisymmetric; a pure rotation needs $a=0$.
\end{proof}

The extension of \cref{prop:EAnotinner} is therefore not one choice among several: at one constituent it is the only equivariant closure. The
frame is abelian, all the dynamics sits in the rotor, and the rotor acts by the binding root itself. With several constituents the answer depends on
the symmetry imposed. If each constituent carries its own copy of $\lG$, so that the symmetry is $\lG^{(1)}\oplus\cdots\oplus\lG^{(n)}$ and no
constituents mix, equivariant maps between distinct constituents vanish, the argument runs factor by factor, and a common rotation rate gives
$\mathfrak s_n=\R\ltimes_{G_n}V_n$ with $G_n=\bigoplus_kG$. Under a common, diagonal Lorentz action this fails: equivalent constituents admit
cross-constituent intertwiners, $\End_{\lG}(V_n)=\{A\otimes\id+B\otimes G\}\cong M_n(\C)$, so the rotor may act by any element of $\uu n$ that
rotates without dilating. When its rates are resonant the frame then admits non-abelian equivariant brackets, already at $n=2$ for rates in the
ratio $1:2$. Under the diagonal action $\R\ltimes_{G_n}V_n$ with an abelian frame is therefore the selected subclass of a common rate and no
mixing, not the only closure.

% --- source: tex/s17_closure#4 (The Lie--Poisson dynamics, and the two interiors as clamped )
\subsection{The Lie--Poisson dynamics, and the two interiors as clamped reductions}
\label{sec:closureLP}

\begin{proposition}[Lie--Poisson equations]
\label{prop:LP}
On $\mathfrak s^*$ with coordinates $(h,p)$ let $H=h^2/(2I_0)+\tfrac12p^{\!\top}Np$, with $I_0>0$ the rotor inertia and $N$ a symmetric form on
$V^*$, the inverse frame inertia. The Lie--Poisson equations are
\begin{equation}
\label{eq:LProtor}
\dot p=\frac h{I_0}\,G^{\!\top}p,\qquad \dot h=-\,p^{\!\top}GNp ,
\end{equation}
so the force law holds with the dynamical rate $\kappa=h/I_0$, and $|p|^2$ is a Casimir. For $p\neq0$ the coadjoint orbit through $(h,p_0)$ is the
cylinder $\R\times\{e^{\theta G}p_0\}\cong T^*S^1$, on which $(\theta,h)$ are canonical for $h^2/(2I_0)+U(\theta)$ with
$U(\theta)=\tfrac12p(\theta)^{\!\top}Np(\theta)$ and $p(\theta)=e^{-\theta G}p_0$: every reduced motion is a motion of the rotor angle alone.
\end{proposition}

\begin{proposition}[The two interiors are clamped reductions]
\label{prop:groupred}
On the group $S$ with product $(\theta,x)(\theta',x')=(\theta+\theta',x+e^{\theta G}x')$ let $L=\tfrac12I_0\dot\theta^2+\tfrac m2\eta(u,u)$ with body velocity
$u=e^{-\theta G}\dot x$, and put $X=e^{-\theta G}x$. Then $u=\dot X+\dot\theta\,GX$, the Euler--Lagrange equations are \eqref{eq:LProtor} with $N=\eta^{-1}/m$,
and clamping $\theta=\kappa\tau$ gives
\[
L=\tfrac12I_0\kappa^2+\tfrac m2\,\eta\bigl(\dot X+\kappa GX,\ \dot X+\kappa GX\bigr),
\]
the split-form interior \eqref{eq:splitL}. With $\tfrac m2|u|^2$ in place of $\tfrac m2\eta(u,u)$ the clamped Lagrangian is the positive-form interior
\eqref{eq:alpha} with $g=G$.
\end{proposition}

So the two interiors of \cref{sec:shifted,sec:neutral} are the clamped reductions of one closure, with inverse inertias $N=\eta^{-1}/m$ and
$N=\id/m$. Which inertias the Clifford data permit is the next question.

% --- source: tex/s17_closure#5 (Two requirements on the inertia)
\subsection{Two requirements on the inertia}

\begin{definition}[The two requirements]
\label{def:inertiareq}
A symmetric $N$ on $V^*$ satisfies the \emph{Lorentz requirement} if it is $\lG$-invariant, $rN+Nr^{\!\top}=0$ for all $r\in\lG$, and the
\emph{duality requirement} if the rate $h/I_0$ is conserved for all data. Write $N=N_c+N_a$ with $N_c=\tfrac12(N-GNG)$ commuting and
$N_a=\tfrac12(N+GNG)$ anticommuting with $G$.
\end{definition}

\begin{theorem}[The two requirements exclude each other]
\label{thm:inertiareq}
\begin{enumerate}[label=(\roman*),leftmargin=2.2em]
\item $N_c$ and $N_a$ are symmetric; the commuting forms have dimension $9$ and the anticommuting ones $12$.
\item The Lorentz requirement gives $N=(a\eta^{-1}+bG\eta^{-1})/m$, a two-dimensional family inside the anticommuting forms, with frame energy
$\tfrac12p^{\!\top}Np=\Re\bigl((a+ib)w(p)\bigr)/(2m)$.
\item The duality requirement holds exactly when $[N,G]=0$. Among the forms invariant under the compact part of $\lG$ the only ones commuting with
$G$ are the multiples of $\id$: the positive form.
\item The two requirements meet in $0$, and they do not exhaust the symmetric forms: $N=(2\,\id+\eta^{-1})/m$ is positive definite and satisfies
neither.
\end{enumerate}
\end{theorem}

In the field-strength reading the dichotomy is the classical one. The Lorentz class consists of the two invariants $|\bB|^2-|\bE|^2$ and
$\bB\cdot\bE$; the duality class contains the energy density $|\bB|^2+|\bE|^2$; and a duality rotation preserves the energy density while it
rotates the Lagrangian density \cite{Calkin65,DeserTeitelboim76,GaillardZumino81}. The split form lies in the Lorentz class and the positive form
in the duality class.

\begin{theorem}[The frame under the closure]
\label{thm:drift}
Reconstruct the group curve from the body velocity,
\[
(\theta,x)^{-1}\tfrac{d}{d\tau}(\theta,x)=(h/I_0,\,Np),\qquad \theta(0)=0 .
\]
Then $p=e^{-\theta G}p_0$ and
\[
\dot x=\bigl(N_c+e^{2\theta G}N_a\bigr)p_0,\qquad \dot h=-\,p_0^{\!\top}Ge^{2\theta G}N_ap_0 .
\]
With the rotor clamped at $\theta=\kappa\tau$, $\kappa\neq0$,
\[
x(\tau)=x_0+\tau\,N_cp_0+\bigl(e^{2\kappa\tau G}-\id\bigr)(2\kappa G)^{-1}N_ap_0 ,
\]
so the frame stays bounded for all data exactly when $N_c=0$ and drifts at the velocity $N_cp_0$ otherwise, while the rate is constant for all data
exactly when $N_a=0$.
\end{theorem}

A clamped rate that is also the dynamical one, and a bounded clamped frame, are therefore exclusive requirements: each is met by one summand of
$\mathrm{Sym}(V^*)=\mathrm{Sym}_c\oplus\mathrm{Sym}_a$. For $N=\id/m$ the rate is conserved and the frame drifts at $p_0/m$, which is the
non-recurrence of \cref{thm:virial}, $X=e^{-\kappa\tau G}(X_0+\tau p_0/m)$; for $N=\eta^{-1}/m$ the clamped frame is bounded, as in \cref{thm:neutral},
but the rate is conserved only on null data.

% --- source: tex/s17_closure#6 (The free rotor: a pendulum in the duality angle)
\subsection{The free rotor: a pendulum in the duality angle}

\begin{theorem}[The duality pendulum]
\label{thm:pendulum}
Let $N=(a\eta^{-1}+bG\eta^{-1})/m$ and write $(a+ib)w(p_0)=R_0e^{i\varphi_0}$, $R_0\ge0$. Along the flow $z(\pi)$ turns as $e^{i\theta}$ and
$w(p)=e^{2i\theta}w(p_0)$, the frame energy is $U(\theta)=\tfrac{R_0}{2m}\cos(2\theta+\varphi_0)$, and
\[
I_0\ddot\theta=\frac{R_0}m\sin(2\theta+\varphi_0),\qquad E_\theta=\tfrac12I_0\dot\theta^2+\frac{R_0}{2m}\cos(2\theta+\varphi_0)=H .
\]
With $\psi:=2\theta+\varphi_0$ this is the pendulum $\ddot\psi=(2R_0/mI_0)\sin\psi$, stable at $\psi=\pi$. The rotor angle is the duality angle of
the momentum field, and the potential is its Lorentz-invariant frame energy.
\end{theorem}

\begin{figure}[t]
\centering
\begin{tikzpicture}
\begin{axis}[width=12cm,height=6.4cm,xmin=-0.3,xmax=6.6,ymin=-3.0,ymax=3.0,
  axis lines=middle,xlabel={$\psi=2\theta+\varphi_0$},ylabel={$\dot\psi$},
  xtick={0,3.14159,6.28318},xticklabels={$0$,$\pi$,$2\pi$},ytick=\empty,
  xlabel style={anchor=north west},ylabel style={anchor=south west}]
\foreach \E/\a in {-0.5/2.0944,0/1.5708,0.5/1.0472}{
  \addplot[tcol,thick,domain=\a:{2*pi-\a},samples=160] {sqrt(max(0,2*(\E-cos(deg(x)))))};
  \addplot[tcol,thick,domain=\a:{2*pi-\a},samples=160] {-sqrt(max(0,2*(\E-cos(deg(x)))))};
}
\addplot[scol,very thick,domain=0:2*pi,samples=200] {2*abs(sin(deg(x/2)))};
\addplot[scol,very thick,domain=0:2*pi,samples=200] {-2*abs(sin(deg(x/2)))};
\addplot[gcol,thick,domain=0:2*pi,samples=160] {sqrt(2*(1.9-cos(deg(x))))};
\addplot[gcol,thick,domain=0:2*pi,samples=160] {-sqrt(2*(1.9-cos(deg(x))))};
\addplot[only marks,mark=*,mark size=2.2pt,black] coordinates {(3.14159,0)};
\addplot[only marks,mark=o,mark size=2.6pt,black] coordinates {(0,0) (6.28318,0)};
\end{axis}
\end{tikzpicture}
\caption{Phase portrait of the duality pendulum (\cref{thm:pendulum}), in units with $2R_0/(mI_0)=1$; the filled dot is the stable point
$\psi=\pi$ and the open dots the unstable one. Blue: librations, in which the rate $\kappa=\dot\theta$ changes sign periodically and the duality
rotation reverses. Red: the separatrix. Green: rotations, in which the rotor turns for all time with a modulated rate. The constant-rate law of \cref{sec:neutral} is the limit of large rotor energy, far above the
separatrix (\cref{cor:multiscale}).}
\label{fig:pendulum}
\end{figure}

\begin{corollary}[When the rate is constant]
\label{cor:constrate}
With Lorentz-class inertia the rate is constant and nonzero along the flow exactly when $R_0=0$ and $h(0)\neq0$; $R_0=0$ means $(a,b)=(0,0)$ or
$w(p_0)=0$, that is $|\bB|=|\bE|$ and $\bB\cdot\bE=0$ for the field of $\pi_0$: a null field, the class that contains plane waves
\cite{Synge56,Robinson61}.
\end{corollary}

\begin{corollary}[The multiscale criterion]
\label{cor:multiscale}
Let $R_0>0$, $\kappa_0:=\dot\theta(0)\neq0$, $\psi_0=2\theta(0)+\varphi_0$ and $\varepsilon:=R_0/(mI_0\kappa_0^2)$. Then
$\kappa^2=\kappa_0^2\bigl[1+\varepsilon(\cos\psi_0-\cos\psi)\bigr]$. The rotor turns for all time exactly when $\varepsilon(1-\cos\psi_0)<1$, and $\varepsilon<\tfrac12$
suffices for every phase, in which case $|\kappa^2/\kappa_0^2-1|\le2\varepsilon$. At $\varepsilon(1-\cos\psi_0)=1$ the motion is the separatrix, and beyond it
the rotor librates and the duality rotation reverses. The constant-rate law is the limit $\varepsilon\to0$, and the multiscale requirement is
\begin{equation}
\label{eq:multiscale}
I_0\,\kappa_0^2\ \gg\ R_0/m .
\end{equation}
\end{corollary}

A clamped rotor is therefore an idealisation of a heavy one. It is not a harmless idealisation, because the frame feels the modulation.

\begin{theorem}[The free rotor drifts the frame]
\label{thm:freedrift}
Let $N$ be of Lorentz class with $R_0>0$ and $\kappa_0\neq0$, in the rotating regime. During one period of the rate, while $\theta$ changes by $\pi$,
the frame is displaced by
\[
\Delta x=\lambda\,e^{-\varphi_0G}Np_0,\qquad \lambda=\tfrac12\int_0^{2\pi}\frac{\cos\psi\,d\psi}{|\kappa_0|\sqrt{1+\varepsilon(\cos\psi_0-\cos\psi)}}>0,
\]
and its mean velocity is $\tfrac\varepsilon4e^{-\varphi_0G}Np_0+O(\varepsilon^2)$. Every rotating solution with $R_0\neq0$ therefore has an unbounded frame.
\end{theorem}

\begin{proof}
For Lorentz-class $N$, $N=N_a$ and $\dot x=e^{2\theta G}Np_0$. In the rotating regime $\dot\theta$ never vanishes, so $d\tau=|d\theta|/|\dot\theta|$, and
with $e^{2\theta G}=e^{-\varphi_0G}e^{\psi G}$ and $e^{\psi G}=\cos\psi+G\sin\psi$ the $\sin\psi$ part integrates to zero because $1/|\dot\theta|$ depends on
$\cos\psi$ only. Since $1/|\dot\theta|$ is a strictly increasing function of $\cos\psi$ when $R_0>0$, the integrand of $\lambda$ exceeds its value at
$\cos\psi=0$ wherever $\cos\psi>0$ and falls short of it where $\cos\psi<0$, so $\lambda>0$. Expanding to first order gives
$\lambda=\pi\varepsilon/(4|\kappa_0|)+O(\varepsilon^2)$ over a period $\pi/|\kappa_0|+O(\varepsilon)$.
\end{proof}

\begin{example}[A drift in closed form]
\label{ex:drift}
Take $p_0=e^1$ and $m=I_0=\kappa_0=1$, so that $w=1$, $\varphi_0=0$, $\varepsilon=1$ and $\dot\theta^2=1+2\sin^2\theta$. Writing
$(1-2s)/\sqrt{1+2s}=2/\sqrt{1+2s}-\sqrt{1+2s}$ with $s=\sin^2\theta$, the displacement per full turn of the rotor, $\theta\mapsto\theta+2\pi$, which is two periods of the rate in
\cref{thm:freedrift}, is
\[
\Delta x=\bigl(8K(-2)-4E(-2)\bigr)e_1\approx0.6336081022\,e_1 ,
\]
with $K(m)=\int_0^{\pi/2}(1-m\sin^2t)^{-1/2}dt$ and $E(m)=\int_0^{\pi/2}(1-m\sin^2t)^{1/2}dt$ the complete elliptic integrals in the parameter $m$
\cite{DLMF}. A numerical integration of the full system, frame included, reproduces this value to ten digits and returns $h$ and $p$ to their
initial values after each turn.
\end{example}

The clamped frame of \cref{thm:drift} moves on a circle of diameter $|Np_0|/|\kappa_0|$, while the free rotor adds $\tfrac{\pi\varepsilon}{2|\kappa_0|}|Np_0|$
per turn to first order, so the constant-rate picture tracks the bounded orbit for about $2/(\pi\varepsilon)$ turns. With Lorentz-class inertia and a
free rotor, a constant rate and a bounded frame coexist exactly on null data.

\begin{proposition}[Several constituents and one rotor]
\label{prop:several}
Let one rotor act on $n$ constituents with $N_k=(a_k\eta^{-1}+b_kG\eta^{-1})/m_k$ and $w_k=w(p^{(k)})$. Then $\dot w_k=2i(h/I_0)w_k$ for every $k$, all
forces carry the common phase $e^{2i\theta}$, and with $R_{\rm eff}e^{i\varphi_{\rm eff}}:=\sum_k(a_k+ib_k)w_k(0)/m_k$ the rotor obeys
$I_0\ddot\theta=R_{\rm eff}\sin(2\theta+\varphi_{\rm eff})$. The \emph{balance surface} $R_{\rm eff}=0$ is preserved by the flow; on it the rate is exactly
constant although no $w_k$ need vanish, and if $h(0)\neq0$ every frame moves on a bounded orbit; with $h(0)=0$ the frames drift linearly. With
equal inertias, $p^{(2)}=-Gp^{(1)}$ gives $w_2=-w_1$.
\end{proposition}

The Clifford data $\eta$, $G$, $G\eta^{-1}$ and $\beta_0$ are integer matrices. The rescaling $(I_0,m)\mapsto(sI_0,tm)$ preserves admissibility and,
at fixed initial rate $\kappa_0$, multiplies $\varepsilon$ by $1/(st)$ (at fixed $h(0)$ it multiplies it by $s/t$), so no value of $\varepsilon$, and no ratio among $I_0$, $m$ and $|a+ib|$, is determined by the Clifford structure:
the form of the inertia is derived, and \eqref{eq:multiscale} is an input.

% --- source: tex/s17_closure#7 (Three constituents: the closure symmetry and the two project)
\subsection{Three constituents: the closure symmetry and the two projectors}
\label{sec:closurethree}

At $n=3$ write $e_A=e^{(k)}_i$ with $A=6(k-1)+i$; the temporal indices are $\{1,2,3,7,8,9,13,14,15\}$, each with its partner $i'=i+3$, and for
temporal $i<j$ put
\[
\mathrm{rot}(i,j):=e_{i'}e_{j'}-e_ie_j,\qquad \mathrm{bst}(i,j):=e_je_{i'}-e_ie_{j'} .
\]
For one constituent these span $\lG$, with $\mathrm{rot}(2,3)=2L_1$ and $\mathrm{bst}(2,3)=2K_1$ and cyclically up to sign. In this numbering the
weak directions of $P_*$ are $W=\{7,13,14,15\}$, the colour directions $\{10,11,12,16,17,18\}$, and $E$ denotes the span of the temporal directions
$\{7,8,9,13,14,15\}$ of constituents two and three; $G_3=G\oplus G\oplus G$.

\begin{theorem}[The closure symmetry is $\so{9,\C}$]
\label{thm:so9C}
$G_3$ is an anti-isometry of $\eta_3$ with $G_3^2=-\id$, and its centraliser in $\so{9,9}$ is $72$-dimensional, spanned by the $36$ paired rotations
and the $36$ paired boosts. In the basis of the nine temporal directions followed by their partners every element is
$\bigl(\begin{smallmatrix}R&S\\-S&R\end{smallmatrix}\bigr)$ with $R,S$ antisymmetric, and $R-iS$ identifies the centraliser with $\so{9,\C}$, whose
compact part consists of the paired rotations.
\end{theorem}

In the field-strength reading $V_3\cong\C^9$ carries three field strengths, and $\so{9,\C}$ consists of the complex orthogonal transformations of
$\sum_kz^{(k)}\cdot z^{(k)}$. The identification of the centraliser of an anti-isometric complex structure with a complex orthogonal algebra is
the inherited linear algebra of Norden geometry \cite{BorowiecFrancavigliaVolovich00}.

\begin{theorem}[No common Lorentz algebra]
\label{thm:nocommon}
The centraliser of $G_3$ meets the interior $\so{1,3}$ of the first constituent in $0$; the elements commuting with both $G_3$ and that algebra form
the $30$-dimensional $\so{6,\C}$ of paired elements of constituents two and three. Under $\lG^{(1)}$ the complexified spinor module is
$128\cdot(\tfrac12,\tfrac12)$, with no Weyl spinor.
\end{theorem}

No nonzero Lorentz generator leaves both the closure and the spinor structure of matter invariant: the two Lorentz algebras are different
subalgebras of $\so{9,9}$ with zero intersection, which is \cref{prop:notconj} lifted to three constituents. A Pati--Salam generator acts on the
closure variables only through its grade-two part, since $\ad_Z$ maps vectors to vectors exactly when $Z$ has grades zero and two. The intersection
with the Pati--Salam algebra can then be computed in general.

\begin{theorem}[The Pati--Salam algebra and the binding root]
\label{thm:PSbinding}
Let $U$ be any four-dimensional subspace of the span of the eight internal temporal directions, and $\mathfrak{ps}_U=\so U\oplus\so6_C$. Then neither
summand meets the centraliser of $G_3$, and
\[
\mathfrak{ps}_U\cap\Cent(G_3)=\bigl\{A+G_3AG_3^{-1}:\ A\in\so{U\cap E}\bigr\}\cong\so{U\cap E},
\]
of dimension $\binom d2\in\{1,3,6\}$ with $d=\dim(U\cap E)$. For the Pati--Salam algebra of $P_*$, $U=\operatorname{span}W\subset E$, and the
binding-compatible part is the diagonal $\so4_\Delta=\operatorname{span}\{\mathrm{rot}(i,j):i<j\in W\}\cong\su2\oplus\su2$, compact and non-abelian.
\end{theorem}

\begin{proof}
Let $A\in\so U$ and $B\in\so6_C$ with $[A+B,G_3]=0$, and split the internal temporal span as $E\oplus F$ with $F=\operatorname{span}\{e_2,e_3\}$. $G_3$ maps
$F$ into the interior plane $\operatorname{span}\{e_5,e_6\}$, where $A$ and $B$ vanish, so $Av=0$ for $v\in F$. For $v\in E$, $G_3v$ is a colour direction
and $Bv=0$, so $G_3Av=BG_3v$ lies in the colour span and $Av\in E$. Hence $A\in\so{U\cap E}$ and $B=G_3AG_3^{-1}$ on the colour span, and conversely.
Finally $d\ge4+6-8=2$.
\end{proof}

Neither Pati--Salam factor contains an element commuting with the binding root; the compatible part pairs every rotated temporal direction with
its spatial partner, which the colour $\su4$, acting on spatial directions alone, never does.

\begin{theorem}[What survives each projector]
\label{thm:survive}
\begin{enumerate}[label=(\roman*),leftmargin=2.2em]
\item The part of $\Cent(G_3)$ commuting with $\Pfam$ is $\operatorname{span}\{\mathrm{rot}(2,3),\mathrm{rot}(8,9),\mathrm{rot}(14,15),\mathrm{bst}(14,15)\}$,
abelian of dimension four. The kinetic condition holds for the two internal rotations and fails for $\mathrm{bst}(14,15)$, so the gauge-admissible
part is $\uu1^2=\operatorname{span}\{\mathrm{rot}(8,9),\mathrm{rot}(14,15)\}$.
\item The part commuting with $P_*$ is $\operatorname{span}\{\mathrm{rot}(2,3),\mathrm{rot}(8,9)\}\oplus\so4_\Delta$, compact of dimension eight, and its
internal part $\so4_\Delta\oplus\R\,\mathrm{rot}(8,9)$ acts faithfully and satisfies the kinetic condition.
\item Under either projector the frame Lorentz algebra $\lG^{(1)}$ keeps only $\mathrm{rot}(2,3)=e_5e_6-e_2e_3$, a rotation of the interior plane
$(e_5,e_6)$ locked to the rotation of the surplus pair.
\end{enumerate}
\end{theorem}

The rotor acts through $G_3$ and commutes with all of $\Cent(G_3)$, so a one-rotor closure is compatible with either projector. What differs is the
symmetry each leaves: an abelian one for $\Pfam$, and for $P_*$ the non-abelian diagonal of weak isospin and a colour $\so4$.

\begin{table}[t]
\centering
\small
\begin{tabular}{@{}lccccccccc@{}}
\toprule
$(q_1,q_2)$ & $(-1,-1)$ & $(-1,0)$ & $(-1,1)$ & $(0,-1)$ & $(0,0)$ & $(0,1)$ & $(1,-1)$ & $(1,0)$ & $(1,1)$\\
\midrule
left, $\Gamma=-1$ & 1 & 0 & 1 & 1 & 2 & 1 & 0 & 2 & 0\\
right, $\Gamma=+1$ & 0 & 2 & 0 & 1 & 2 & 1 & 1 & 0 & 1\\
\bottomrule
\end{tabular}
\caption{Weyl content of the block of $\Pfam$ under the two surviving charges $Q_1=-\tfrac12J\,\mathrm{rot}(8,9)$ and $Q_2=-\tfrac12J\,\mathrm{rot}(14,15)$:
the number of Weyl components with charges $(q_1,q_2)$ on each chirality half.}
\label{tab:closurecontent}
\end{table}

\begin{theorem}[Chirality and anomalies of what survives]
\label{thm:closurechiral}
With anomaly coefficients $\mathcal A(f)=\tfrac14\bigl[\Tr(fP_L)-\Tr(fP_R)\bigr]$ and charges $-JB$:
\begin{enumerate}[label=(\roman*),leftmargin=2.2em]
\item On $\Pfam$ the content is \cref{tab:closurecontent}; $\mathcal A\bigl((\alpha Q_1+\beta Q_2)^3\bigr)=-12\alpha\beta^2$ and $\mathcal A(\alpha Q_1+\beta Q_2)=0$, and
the left and right charge multisets coincide exactly when $\alpha\beta=0$. The anomaly-free subalgebras of $\uu1^2$ are $0$, $\R Q_1$ and $\R Q_2$,
each vector-like.
\item On $P_*$, with $\so4_\Delta=\su2_+\oplus\su2_-$ and $Q=-\tfrac12J\,\mathrm{rot}(8,9)$, the halves are
$L=(0,0)_0\oplus(\tfrac12,\tfrac12)_{+1}\oplus(1,0)_0$ and $R=(0,0)_0\oplus(\tfrac12,\tfrac12)_{-1}\oplus(0,1)_0$. For $H=\alpha H^+_3+\beta H^-_3+\gamma Q$,
$\mathcal A(H^3)=2\gamma(3\alpha^2+3\beta^2+4\gamma^2)$ and $\mathcal A(H)=8\gamma$, so the anomaly-free subalgebras are exactly those of $\so4_\Delta$, which
has neither a perturbative nor a global anomaly.
\item The two halves of $P_*$ carry inequivalent representations of $\so4_\Delta$, so it acts chirally; but every summand is a real representation
of $\su2\oplus\su2$, so the left-handed content is real and gauge-invariant Majorana masses exist \cite{Suzuki2000}. The chirality of $\so4_\Delta$
on $P_*$ forbids no mass.
\end{enumerate}
\end{theorem}

The rotor has no continuous lift to spinors: $(e^{\theta G})^{\!\top}\eta\,e^{\theta G}=\cos2\theta\,\eta+\sin2\theta\,\eta G$, so $e^{\theta G}$ is an isometry of
$\eta$ only for $\theta\in\pi\Z$, and at $\theta=\pi$, $e^{\pi G_3}=-\id$ is implemented by $\Gamma$. This excludes a coupling of the rotor through the
spin representation and says nothing about other couplings.

\begin{theorem}[Inertia freedom at three constituents]
\label{thm:inertia3}
The real dimensions of the symmetric forms on $V_3^*$ in the Lorentz class, $\{N,G_3\}=0$, and in the duality class, $[N,G_3]=0$, that are invariant
under a given symmetry are:
\begin{center}
\begin{tabular}{@{}lcc@{}}
\toprule
symmetry & Lorentz class & duality class\\
\midrule
$\Cent(G_3)=\so{9,\C}$ & $2$ & $0$\\
$\bigoplus_k\lG^{(k)}$ & $6$ & $0$\\
the $\Pfam$-surviving algebra & $18$ & $13$\\
the $P_*$-surviving algebra & $8$ & $6$\\
none & $90$ & $81$\\
\bottomrule
\end{tabular}
\end{center}
\end{theorem}

The Lorentz-class inertia is fixed up to one complex scale under the full closure symmetry and up to three under the single-constituent Lorentz
algebras, and imposing only what a projector leaves enlarges the freedom to nine complex scales for $\Pfam$ and four for $P_*$: neither projector
selects the inertia.

% --- source: tex/s17_closure#8 (Brackets supplied by hand)
\subsection{Brackets supplied by hand}

The closure above is derived; one may also supply a bracket on the frame by hand and ask what the arena's own inertias do with it. A simple
quadratic Lie algebra whose invariant form is $\eta$ must have a maximal compact subalgebra of exactly half its dimension. For algebras that
are not simple this fails: $\so4$ with invariant forms of opposite sign on its two factors is compact and has signature $(3,3)$. Such
\emph{balanced} simple real forms are rare: a complex simple algebra regarded as real, $\su{p,q}$ with $|p-q|=1$, or $\so{p,q}$ with $(p,q)=\bigl(\tfrac{k(k+1)}2,\tfrac{k(k-1)}2\bigr)$, and
no exceptional form. At the first rung $\R^{3,3}$ admits exactly three semisimple quadratic structures, $\so{1,3}$, $\so4$ and $\so{2,2}$, and only
$\so{1,3}$ if $\eta$ is to be proportional to the Killing form. For $n\le20$ a simple closure exists exactly at $n=1,4,5,6,7,8,12,15,16,20$.

Two observations about the arena's inertias hold across these closures. With the split-form inertia $A=m\eta$, which is an invariant form, every
solution of the Euler--Arnold equation is a one-parameter subgroup and $\dot p\equiv0$. With the positive form $A=m\,\id$ on $\so{1,3}$ the system is
$\dot\bB=0$, $\dot\bE=\tfrac2m\bE\times\bB$: the magnetic half is frozen and the electric half precesses about it. The inertias with
$G^{\!\top}A^{-1}G=-A^{-1}$ form a twelve-dimensional family of holomorphic Euler tops $\dot Z=Z\times MZ$ with $M$ complex symmetric; they are Liouville
integrable and never positive definite.

At the rungs $n=(N^2-1)/3$ with $3\nmid N$, that is $(N,n)=(2,1),(4,5),(5,8),\dots$, the compact half of the balanced form $\slr{N,\C}_\R$ is $\su N$,
which is Zeitlin's finite truncation of two-dimensional incompressible Euler flow \cite{Zeitlin91}. Giving the 't~Hooft clock--shift modes the
discrete-Laplacian inertia, numerical integration at $N=4$ and $N=5$ conserves energy and enstrophy to $10^{-12}$ and moves up to $37\%$ of the
energy and $54\%$ of the enstrophy between Laplacian shells: transfer between scales inside the frame. With either of the arena's own inertias the same
data are exact equilibria, because both $\id$ and $\eta$ restrict to multiples of the invariant form on $\su N$. What the frame lacks for a fluid is
the inertia. Finally, no closure dissipates: every Euler--Arnold flow stays on a coadjoint orbit and preserves every Casimir, which is exactly what
viscosity destroys, so no closure of the force law can produce $\nu\Delta u$ on its own; and a constant generator that uniformly decreases the split
form $C$ has, by the inertia theorem of Ostrowski and Schneider \cite{OstrowskiSchneider62}, exactly $3n$ exponentially unstable directions.

% --- source: tex/s17b_rest#0 ((head))


% --- source: tex/s17b_rest#1 (Rest energy in the closure)
\section{Rest energy in the closure}
\label{sec:restclosure}

The closure of \cref{sec:closure} is the one dynamics that the arena closes by itself: the rate $\kappa$ of the shift operators
$\alpha_\pm=\partial_\tau\pm\kappa G$ becomes the angular velocity of a rotor. This section asks what that dynamics does to the energies of a
field point, the undivided energy $\tfrac12(|\bB|^2+|\bE|^2)$ and the rest energy, and whether it relates the rest energy to the rate. Two fields
are kept apart throughout: the \emph{momentum field} $\pi=\eta^{-1}p$, which is $m\alpha_+X$ in the clamped split-form interior, and the
\emph{frame field} $X$. The answers, in brief, are these. The closure freezes the undivided energy, the rest energy and the Poynting vector of the momentum
field, for every Hamiltonian; it changes those of the frame field only by an oscillation or an unbounded drift, and selects no value. The rest
energy of the momentum field is exactly the amplitude of the potential of the duality angle, and the conserved total splits into a rate part and a
Lorentz-invariant frame part, whose exchange of energy is fixed by that rest energy alone. The constant rate of the interior is the limit of a heavy rotor. The rest flow of Dirac
matter and the binding rotation of the frame are two rotations at rest that the spin homomorphism does not relate; $m$ and $\kappa$ are independent inputs, and $mc^2=\hbar\kappa$ is an added
identification, and nothing in the algebra fixes the length $\ell$ in $\kappa=c/\ell$. A quantised rotor gives a rate ladder only under stated
assumptions.

\paragraph{Notation.} As in \cref{sec:closure}, $V=\R^{3,3}$, $G$ is the binding root, $G^2=-\id$, $G^{\!\top}=-G$, $\eta G=-G\eta$, a covector
$p$ is a column, and $z(Y)=\bB+i\bE$ for $Y\in V$ with components $(\bB;\bE)$ along $(e_1,e_2,e_3;e_4,e_5,e_6)$, so that $G$ acts as $i$
(\cref{thm:emfield}). With $q(z)=z\cdot z$, the undivided energy $E_6(Y)$, the rest energy $E_0(Y)$ and the Poynting vector
$\mathbf S(Y)=\bE\times\bB$ of $Y$, minus the textbook one, are those of \eqref{en:defs}; only $|\mathbf S|$ enters below. For the momentum field, $E_6(\pi)=\tfrac12|p|^2$ because $\eta$
is orthogonal, and $E_0(\pi)=\tfrac12|w(p)|$ with $w$ as in \eqref{eq:wdef}. The Lorentz-class inertia of \cref{thm:inertiareq}(ii) is written
$N=(\nu_1\eta^{-1}+\nu_2G\eta^{-1})/m$, $\nu:=\nu_1+i\nu_2$, its coefficients being the $(a,b)$ of that theorem. Units are $\hbar=c=1$ unless a
symbol is shown, $\tau$ is a time, and $\mathrm K(\mu)$, $\mathrm E(\mu)$ are the complete elliptic integrals of \cref{ex:drift}.

% --- source: tex/s17b_rest#2 (What the closure conserves)
\subsection{What the closure conserves}
\label{sec:restconserve}

On $\mathfrak s^*$, the dual of the unique closure $\mathfrak s=\R\hat g\ltimes_GV$ of \cref{thm:closureunique}, with coordinates $(h,p)$,
$F_h=\partial F/\partial h$ and $\nabla F=\partial F/\partial p$, the Lie--Poisson bracket and the equations of an arbitrary Hamiltonian $H$ are
\begin{equation}
\label{eq:restLPgen}
\{F_1,F_2\}=p^{\!\top}\bigl(F_{2,h}\,G\nabla F_1-F_{1,h}\,G\nabla F_2\bigr),\qquad
\dot p=H_h\,G^{\!\top}p,\qquad \dot h=-\,p^{\!\top}G\,\nabla H .
\end{equation}
For $H=h^2/(2I_0)+\tfrac12p^{\!\top}Np$ these are the equations \eqref{eq:LProtor} of \cref{prop:LP}.

\begin{theorem}[The momentum field is frozen up to its duality phase]
\label{cl:frozen}
\begin{enumerate}[label=(\roman*),leftmargin=2.2em]
\item A function $f(p)$ is a Casimir of $\mathfrak s^*$ if and only if it is invariant under $p\mapsto e^{tG}p$, that is under the duality rotations
$z(\pi)\mapsto e^{-it}z(\pi)$ of the momentum field. Hence $E_6(\pi)$, $E_0(\pi)^2$ and the vector $\mathbf S(\pi)$ are smooth Casimirs, while $\Re w$
and $\Im w$ are not; $E_0=|w|/2$ is conserved as well but is not smooth at $w=0$.
\item At $p\neq0$ the symplectic leaves are two-dimensional, so there are five independent local Casimirs. The real and imaginary parts of the
products $\bar z_jz_k$, $z=z(\pi)$, provide five, and $E_6(\pi)$, $E_0(\pi)^2$, $\mathbf S(\pi)$ account generically for four of them, subject to
$E_6(\pi)^2=E_0(\pi)^2+|\mathbf S(\pi)|^2$.
\item Along the flow of every Hamiltonian, $p(\tau)=e^{-\theta(\tau)G}p_0$ with $\dot\theta=H_h$ and $\theta(0)=0$; so
$z(\pi(\tau))=e^{i\theta(\tau)}z(\pi_0)$ and $w(p(\tau))=e^{2i\theta(\tau)}w(p_0)$.
\end{enumerate}
\end{theorem}

\begin{proof}
(i) For $f=f(p)$, \eqref{eq:restLPgen} gives $\{f,F\}=F_h\,p^{\!\top}G\nabla f=-F_h\,(Gp)\cdot\nabla f$, which vanishes for every $F$ exactly when $f$
is constant along $t\mapsto e^{tG}p$. Since $\eta G=-G\eta$, $\eta^{-1}e^{tG}p=e^{-tG}\pi$, on which $G$ acts as $i$. The quantities $|z|^2$,
$|z\cdot z|$ and $\bar z\times z$ are invariant under a phase, while $w=z\cdot z$ is multiplied by $e^{-2it}$. (ii) From \eqref{eq:restLPgen},
$\{p_i,p_j\}=0$ and $\{h,p_j\}=(Gp)_j$, so the Poisson tensor has rank two at $p\neq0$ and the leaves have codimension five in the
seven-dimensional $\mathfrak s^*$. The map $z\mapsto z\bar z^{\!\top}$ from $\C^3\setminus\{0\}$ onto the rank-one positive Hermitian matrices has the
phase circles as its fibres, so its image is five-dimensional and its nine real components contain five independent functions, all phase
invariant. The relation is the identity \eqref{en:lagrange} of \cref{en:split}. (iii) $\dot p=-H_hGp$ integrates to
$p=e^{-\theta G}p_0$, so $\pi=e^{\theta G}\pi_0$, and $w$ is quadratic in $z$.
\end{proof}

The theorem is fixed by the algebra. In substance it is a corollary of \cref{prop:LP}, whose coadjoint orbits are the cylinders
$\R\times\{e^{\theta G}p_0\}$; item (iii) is the phase law of \cref{thm:pendulum}, extended to every Hamiltonian, and the relation in (ii) is
inherited. That $E_6(\pi)$, $E_0(\pi)$ and $\mathbf S(\pi)$ are conserved for every Hamiltonian follows from \cref{prop:LP}, and that they are constant already in the clamped
interior is \cref{en:interior}(i). What the theorem adds is that the Casimirs are exactly the duality-invariant functions of $p$, and their count.

\begin{corollary}[What the closure can and cannot change]
\label{cl:nogen}
\begin{enumerate}[label=(\roman*),leftmargin=2.2em]
\item For every Hamiltonian on $\mathfrak s^*$, quadratic or not, of Lorentz class or not, $E_6(\pi)$, $E_0(\pi)$ and $\mathbf S(\pi)$ are constant
along every solution: a null momentum field stays null, and a momentum field with rest energy $E_0(\pi_0)$ keeps it.
\item The same functions of the frame field are not constants of motion. In the clamped split-form interior, where
$z(X)=c_-e^{-i\kappa\tau}+c_+e^{i\kappa\tau}$, they oscillate in general (\cref{en:interior}(ii)); for instance $c_-=(1,i,0)$, $c_+=(1,0,0)$ give
$|q(z(X))|^2=9-8\sin^2\kappa\tau$. Under the
free-rotor drift of \cref{ex:drift} with $x(0)=0$, the frame field $X=e^{-\theta G}x$ has $E_0(X)=0$ at $\tau=0$ and
$E_0(X)=E_6(X)=\tfrac12k^2|\Delta x|^2$ after $k$ full turns, $\Delta x$ the displacement per turn of that example; after three turns this is
$\approx1.807$. On every rotating solution with Lorentz-class inertia and $R_0>0$, in the notation of \cref{thm:pendulum}, $E_6(X)$ and
$E_0(X)$ are unbounded.
\end{enumerate}
\end{corollary}

\begin{proof}
(i) Casimirs Poisson-commute with every Hamiltonian (\cref{cl:frozen}(i)). (ii) By the expansion of $z\cdot z$ in \cref{en:interior}(ii), in the example $c_-\cdot c_-=0$ and $c_-\cdot c_+=c_+\cdot c_+=1$, so $|q|^2=|2+e^{2i\kappa\tau}|^2=5+4\cos2\kappa\tau$. Since $G$ acts as $i$,
$e^{-\theta G}$ multiplies $z$ by a phase, so $E_6(X)=E_6(x)$ and $E_0(X)=E_0(x)$. In \cref{ex:drift} the motion of $(\theta,h)$ and the velocity
$\dot x=e^{2\theta G}Np_0$ are periodic in $\theta$, so at $\theta=2\pi k$ the frame is $x=k\,\Delta x$, a real multiple of $e_1$, and
$q(z(x))=k^2|\Delta x|^2$. In general, by \cref{thm:freedrift} the frame advances by $\Delta x=\lambda e^{-\varphi_0G}Np_0$, $\lambda>0$, in each
period of the rate, so $x=x_0+n\Delta x$ after $n$ periods. As $z(Np_0)=\nu\,z(\pi_0)/m$, $q(z(\Delta x))=\lambda^2e^{-2i\varphi_0}\nu^2w(p_0)/m^2$,
which is nonzero because $|\nu w(p_0)|=R_0>0$; hence $|q(z(x))|$ and $|x|^2$ grow like $n^2$.
\end{proof}

Item (i) is an obstruction, under the hypothesis that the dynamics is a Hamiltonian flow on $\mathfrak s^*$: no such flow produces, removes or
changes the rest energy of the momentum field. Item (ii) is fixed: its statement on the clamped interior is \cref{en:interior}(ii), its growth follows from \cref{prop:orbitsEM,thm:freedrift}, and the two examples are added. The first verdict therefore cuts both
ways. For the momentum field, $E_6(\pi)$ and its split $E_6(\pi)^2=E_0(\pi)^2+|\mathbf S(\pi)|^2$ label the leaf: the rest energy is carried by the
dynamics, not produced by it. For the frame field, to which the field-strength reading of $E_6$ and $E_0$ refers, the closure does change the rest
energy, by an oscillation or by an unbounded drift, but not by any mechanism that selects a value. Nor is $E_6(X)$ a function of
$E_6(\pi)$: on the orbits without momentum with $X\neq0$, $|p|^2=0<E_6(X)$ (\cref{en:interior}(iii)). What moves is the duality angle $\theta$
and the rate $h/I_0$.

% --- source: tex/s17b_rest#3 (The duality pendulum)
\subsection{The duality pendulum}
\label{sec:restpendulum}

Unless duality-class inertia is named, $N$ is from now on of Lorentz class, $\nu w(p_0)=R_0e^{i\varphi_0}$ with $R_0\ge0$, and $\psi:=2\theta+\varphi_0$. By \cref{thm:inertiareq}(ii) the
frame energy is $\tfrac12p^{\!\top}Np=\Re(\nu w(p))/(2m)$, and by \cref{thm:pendulum} the rotor is the pendulum
\[
\frac{h^2}{2I_0}+U(\theta),\qquad U(\theta)=\frac{R_0}{2m}\cos\psi,\qquad I_0\ddot\theta=\frac{R_0}m\sin\psi,\qquad
R_0=|\nu|\,|w(p_0)|=2|\nu|\,E_0(\pi_0).
\]
The amplitude of the potential of the duality angle is therefore the rest energy of the momentum field, times $|\nu|/m$; its undivided energy
$E_6(\pi)$ and its Poynting vector $\mathbf S(\pi)$ do not enter.

\begin{proposition}[Small oscillations]
\label{cl:omega}
The stable equilibrium is $\psi=\pi$, where $U=-R_0/(2m)$. With $\psi=\pi+\xi$, $\ddot\xi=-\Omega^2\sin\xi$ and
\[
\Omega^2=\frac{2R_0}{mI_0}=\frac{4|\nu|\,E_0(\pi_0)}{mI_0} ;
\]
$\theta$ oscillates at the same frequency with half the amplitude of $\xi$, and at amplitude $0<\xi_0<\pi$ the period is
$4\,\mathrm K\bigl(\sin^2(\xi_0/2)\bigr)/\Omega$.
\end{proposition}

\begin{proof}
$\ddot\psi=2\ddot\theta=(2R_0/mI_0)\sin\psi$ and $\sin(\pi+\xi)=-\sin\xi$. The period follows from $\tfrac12\dot\xi^2=\Omega^2(\cos\xi-\cos\xi_0)$ with
the substitution $\sin(\xi/2)=\sin(\xi_0/2)\sin t$.
\end{proof}

The proposition is fixed, as a corollary of \cref{thm:pendulum}, which it linearises; the period is the inherited one of the pendulum
\cite{DLMF}. By \cref{cor:constrate}, with Lorentz-class inertia the rate is constant and nonzero exactly when $R_0=0$, that is when $\nu=0$ or
the momentum field is null, and $h(0)\neq0$; then $U\equiv0$ and $\theta=h\tau/I_0$ is an exact flat direction. For duality-class inertia, $[N,G]=0$, the rate is constant for all data
(\cref{thm:inertiareq}(iii), \cref{thm:drift}), and the potential $U(\theta)=\tfrac12p_0^{\!\top}e^{\theta G}Ne^{-\theta G}p_0=\tfrac12p_0^{\!\top}Np_0$ is
constant, in general not zero. The remaining case is the following.

\begin{corollary}[A constant rate with a rest energy]
\label{cl:flat}
Let $N$ be of Lorentz class with $\nu\neq0$, and let $E_0(\pi_0)>0$. The rate $h/I_0$ is constant along a solution if and only if $h(0)=0$ and
$\psi(0)\in\pi\Z$. There the rate vanishes and the frame moves at the constant velocity $Np_0$, with $|Np_0|=|\nu|\,|p_0|/m\neq0$.
\end{corollary}

\begin{proof}
$\dot h=(R_0/m)\sin\psi$ with $R_0>0$. If $h$ is constant, $\sin\psi\equiv0$, so $\theta$ is constant and $h=I_0\dot\theta=0$; conversely these points are
equilibria. With $\theta\equiv0$, \cref{thm:drift} ($N_c=0$) gives $\dot x=Np_0$, and $N^2=|\nu|^2\id/m^2$ because $\eta^{-1}G\eta^{-1}=-G$.
\end{proof}

The corollary is fixed by the algebra. It complements \cref{cor:constrate}, which covers zero rest energy, and adds the uniform motion of the frame
at the equilibria.

\begin{remark}[A pseudo-Goldstone reading, and its limits]
\label{cl:goldstone}
On a leaf, $\theta$ is the angle of the duality group $U(1)_G$ of the momentum field, which a nonzero rest energy breaks to $\{\pm\id\}$ on the
space of Lorentz orbits (\cref{en:doubly}; \cref{sec:massverdict}). In the duality class it is a symmetry of $H$; in the Lorentz class it is broken on each leaf in proportion to
$E_0(\pi)$, and exactly
\[
I_0\,\Omega^2=U''(\theta_{\min})=\frac{2R_0}m=\frac{4|\nu|\,E_0(\pi_0)}m .
\]
The first equality defines $\Omega$ for every oscillator with one degree of freedom; the content specific to the closure is the second, which is
\cref{thm:pendulum}. With $I_0$ in the place of the square $f_\pi^2$ of the pion decay constant and $|\nu|E_0(\pi_0)/m$ in the place of the explicit breaking, the relation has the
shape of the Gell-Mann--Oakes--Renner relation for a pseudo-Goldstone boson \cite{GellMannOakesRenner68,Weinberg72}, as every pendulum does.
That is a reading, and its limits are definite. The breaking is explicit, by the inertia, not spontaneous: there is no vacuum manifold, and
$E_0(\pi)$ is a label fixed by the data (\cref{cl:nogen}(i)), not an order parameter; and by the bound \eqref{eq:EBbound} the energy of a field point
is least at the zero field, with no degenerate family of lowest-energy configurations among which a phase of $q$ could be chosen. With one degree of freedom, ``mass'' here means a frequency. What survives the reading is exact: on a null
momentum field the angle is a zero mode, and a rest energy of the momentum field gives it the frequency $\Omega\propto\sqrt{E_0(\pi_0)}$.
\end{remark}

% --- source: tex/s17b_rest#4 (The energy budget: a rate part and a frame part)
\subsection{The energy budget: a rate part and a frame part}
\label{sec:restbudget}

\begin{proposition}[Rate part and frame part]
\label{cl:budget}
With Lorentz-class inertia, along every solution:
\begin{enumerate}[label=(\roman*),leftmargin=2.2em]
\item $H=K_{\rm rot}+U$ with $K_{\rm rot}=\tfrac12I_0\kappa^2\ge0$, $\kappa=h/I_0$, $U=\tfrac{R_0}{2m}\cos\psi=\Re(\nu\,w(p))/(2m)$, and
$\dot K_{\rm rot}=-\dot U=\tfrac{R_0}m\,\kappa\sin\psi$;
\item in the rotating regime $K_{\rm rot}$ runs through exactly $\bigl[H-\tfrac{R_0}{2m},\,H+\tfrac{R_0}{2m}\bigr]$ in every period of the rate, so the
energy exchanged between the two parts is $R_0/m=2|\nu|E_0(\pi_0)/m$; in the librating regime it runs through $\bigl[0,\,H+\tfrac{R_0}{2m}\bigr]$;
\item $K_{\rm rot}$ and $U$ are separately invariant under $\lG$ acting by $(h,p)\mapsto(h,\Lambda^{-\top}p)$, $\Lambda=e^r$, $r\in\lG$, and the flow
commutes with this action;
\item $(\theta,h)(\tau)$ depends on $p_0$ only through $\nu w(p_0)$; $E_6(\pi_0)$ and $\mathbf S(\pi_0)$ do not enter the motion.
\end{enumerate}
\end{proposition}

\begin{proof}
(i) $\dot K_{\rm rot}=h\dot h/I_0$ and $\dot U=-(R_0/2m)\,2\kappa\sin\psi$, with $\dot h=(R_0/m)\sin\psi$. (ii) In rotation $\dot\psi=2\kappa$ never
vanishes and $\cos\psi$ covers $[-1,1]$ in each period; in libration $K_{\rm rot}=0$ at the turning points and is largest at $\psi=\pi$. (iii) The
Lorentz requirement $rN+Nr^{\!\top}=0$ of \cref{def:inertiareq} integrates to $\Lambda^{-1}N\Lambda^{-\top}=N$, so $U(\Lambda^{-\top}p)=U(p)$; since
$[r,G]=0$, $\Lambda^{-\top}$ commutes with $G^{\!\top}$, and both equations \eqref{eq:LProtor} are preserved. (iv) $\eta^{-1}\Lambda^{-\top}=\Lambda\eta^{-1}$
because $\Lambda\in O(3,3)$, and $\lG\cong\so{3,\C}$ acts on $z$ complex-linearly and preserves $q$ (\cref{thm:GLorentz}), so
$w(\Lambda^{-\top}p)=w(p)$; the equations for $(\theta,h)$ contain only $R_0$ and $\varphi_0$.
\end{proof}

Item (i) restates the energy $E_\theta=H$ of \cref{thm:pendulum}; items (ii)--(iv) are fixed by the algebra. The conserved total thus
differentiates into two Lorentz-invariant parts, a rate part and a frame part, whose exchange is fixed by the rest energy of the momentum field
alone; neither part is the undivided energy $E_6(\pi)$, which is inert.

\paragraph{Units, and the factor one half.} Take $[\tau]=\mathsf T$ and $[\kappa]=\mathsf T^{-1}$, as in \eqref{eq:kappa}, and $[X]=\mathsf L$, a
convention by \cref{cl:scalings}(i) under which $m$ is a mass. Then $[I_0]=\mathsf{ML}^2$, $[h]=\mathsf{ML}^2\mathsf T^{-1}$ is an action,
$[p]=\mathsf{MLT}^{-1}$, $[N]=\mathsf M^{-1}$ and $[H]=[K_{\rm rot}]=[U]=\mathsf{ML}^2\mathsf T^{-2}$; $E_0(\pi)$ is a squared momentum and
$|\nu|E_0(\pi)/m$ an energy. With the clamped value $\kappa=c/\ell$ and $M_{\rm rot}:=I_0/\ell^2$, a mass,
\[
K_{\rm rot}=\tfrac12I_0\kappa^2=\tfrac12M_{\rm rot}c^2 .
\]
Reading this as the ``$c^2$ part'' of the conserved total is a reading: it is literally the Newtonian kinetic energy of a mass $M_{\rm rot}$ on a
circle of radius $\ell$ at rim speed $c$, outside the domain of Newtonian kinetics. It is not $E=mc^2$. The factor $\tfrac12$ comes from an energy
quadratic in the rate, where Planck's relation is linear; $M_{\rm rot}$ is not tied to $m$, since the rescaling $(I_0,m)\mapsto(sI_0,tm)$ recalled in \cref{sec:restscale} changes the two
independently; and $\kappa=c/\ell$ holds only
for the clamped rotor.

% --- source: tex/s17b_rest#5 (When the rate is constant: the multiscale criterion)
\subsection{When the rate is constant: the multiscale criterion}
\label{sec:restmultiscale}

\Cref{cor:multiscale} answers when the rate is nearly constant. In terms of the two parts of the energy its parameter is
\[
\varepsilon=\frac{R_0}{mI_0\kappa_0^2}=\frac{U_{\max}}{K_{\rm rot}(0)},\qquad U_{\max}=\frac{R_0}{2m}=\frac{|\nu|\,E_0(\pi_0)}m ,
\]
the ratio of the rest-energy amplitude of the momentum field to the initial rotor energy. If $c$ is the fixed conversion of \cref{thm:onec}, a
varying rate $\kappa=c/\ell$ is a varying length $\ell(\tau)=c/\kappa(\tau)$, and the law $\kappa^2=\kappa_0^2[1+\varepsilon(\cos\psi_0-\cos\psi)]$ of
\cref{cor:multiscale} gives it the relative spread $\varepsilon+O(\varepsilon^2)$ along a rotating solution. With $\kappa_0=c/\ell$, the multiscale
requirement \eqref{eq:multiscale} becomes
\[
\tfrac12M_{\rm rot}c^2\ \gg\ \frac{|\nu|\,E_0(\pi_0)}m :
\]
for ``$c/\ell$'' to be constant, the rotor energy must dominate the rest-energy amplitude of the momentum field. The clamped split-form interior
of \cref{sec:neutral} is the limit $\varepsilon\to0$ (\cref{prop:groupred,cor:multiscale}); the positive-form interior of \cref{sec:shifted} has
duality-class inertia and an exactly constant rate (\cref{thm:inertiareq}(iii)). Read in this way, the shift $\kappa G$, with $(\kappa G)^2=-c^2/\ell^2\,\id$, which \cref{sec:energy-readings} reads as the ``symmetric part'' of $E=mc^2$
(\cref{en:onec}(iii), \cref{thm:EAdichotomy}), is constant only insofar as the differentiated part is negligible beside the rotor; $\varepsilon$ is an input, and the
reading is ours.

% --- source: tex/s17b_rest#6 (Two rotations at rest)
\subsection{Two rotations at rest}
\label{sec:resttworot}

Two rotations at rest are available. The first belongs to Dirac matter. At $\bp=0$ the Hamiltonian \eqref{eq:Hp} of the commuting
first-order form $\cL_1=\half\Psi^{\!\top}J\dot\Psi-\half\Psi^{\!\top}H\Psi$ is $H(0)=m\gamma^0=me_1$.

\begin{proposition}[The rest flow]
\label{mb:restflow}
$JH(0)=mT$, so the rest flow is $\Psi(\lambda)=e^{-m\lambda T}\Psi_0=(\cos m\lambda\,\id-\sin m\lambda\,T)\Psi_0$. On the eigenspaces $e_1=\pm1$, each
four-dimensional and $J$-invariant, $T=\pm J$, so $\Psi(\lambda)=e^{\mp m\lambda J}\Psi_0$ with energy $\half\Psi^{\!\top}H\Psi=\pm\half m|\Psi|^2$.
Moreover $\gamma^0=-TJ$.
\end{proposition}

\begin{proof}
$Je_1=e_2e_3e_1=e_1e_2e_3=T$. Since $T=e_1J$ with $[e_1,J]=0$, $T\Pi_\pm=\pm J\Pi_\pm$ for $\Pi_\pm=\half(\id\pm e_1)$; $T^2=-\id$ gives the
exponential, and $-TJ=-e_1J^2=e_1$.
\end{proof}

The proposition is inherited: it is the realification of $\psi=e^{\mp im\lambda}\psi_0$ on $\beta=\pm1$, at angular frequency $mc^2/\hbar$
in restored units, de Broglie's internal clock \cite{deBroglie1925}. The second rotation belongs to the frame: the
orbits without momentum of the split-form interior are the duality rotations $z\mapsto e^{-i\kappa\tau}z$ of angular frequency $\kappa=c/\ell$,
along which $q$ turns at $2\kappa$, and the duality rotations of the other sense carry all the momentum (\cref{prop:orbitsEM}(i)--(ii),
\cref{en:interior}(iii)); the sense without momentum is singled out by the momentum, not by duality.

\begin{proposition}[No operator relates the two rates]
\label{cl:norel}
\begin{enumerate}[label=(\roman*),leftmargin=2.2em]
\item $\Ad_{e^{\vartheta T}}$ fixes $e_1,e_2,e_3$ and maps a spatial generator $e_a$ to $\cos2\vartheta\,e_a+\sin2\vartheta\,Te_a$, which is a vector
only for $\vartheta\in\tfrac\pi2\Z$. The Dirac rest flow is therefore not a motion of $V$, and its generator $T$, of grade three, is not in $\so{3,3}$.
\item $(e^{\theta G})^{\!\top}\eta\,e^{\theta G}=\cos2\theta\,\eta+\sin2\theta\,\eta G$, so $e^{\theta G}$ lies in $O(3,3)$, and lifts to
$\mathrm{Pin}(3,3)$, only for $\theta\in\pi\Z$.
\end{enumerate}
Hence neither one-parameter group is the image of the other under the spin homomorphism, and this route supplies no relation between $m$
and $\kappa$.
\end{proposition}

\begin{proof}
(i) $T$ commutes with the temporal and anticommutes with the spatial generators (\cref{lem:blade}), so $e^{\vartheta T}e_ae^{-\vartheta T}=e^{2\vartheta T}e_a$,
and $Te_a$ has grade four. (ii) Expand with $e^{\theta G}=\cos\theta+\sin\theta\,G$, using that $G^{\!\top}\eta=\eta G$ is symmetric and
$G^{\!\top}\eta G=-\eta$.
\end{proof}

The proposition is an obstruction to one route only, a map of the arena carrying one rate to the other; item (ii) is the remark closing the
three-constituent part of \cref{sec:closure}, and the corresponding statement in \cref{sec:orth}. That $m$ and $\kappa$ are independent is
not a consequence of it: both are free inputs. What remains is an identification, and it is added, not derived. With $E=\hbar\times$rate,
$\hbar\kappa=mc^2$ holds if and only if $\ell=\hbar/(mc)$, the reduced Compton wavelength. The dynamics does not supply it. The interior equation
$\ddot X=-\kappa^2X$ does not contain $m$, which multiplies the Lagrangian; and the rescaling $(\tau,\kappa,I_0,m)\mapsto(s\tau,\kappa/s,sI_0,sm)$ of
\cref{cl:scalings}, which maps solutions to solutions and preserves the action, hence any quantisation of $h$, keeps $m\kappa$ and multiplies
$m/\kappa$ by $s^2$, so that $m=\hbar\kappa/c^2$ picks one point of each orbit by fiat. The rate is moreover ambiguous by a factor two: $q$ turns at
$2\kappa$ while $z$ turns at $\kappa$, and matching $2\hbar\kappa=mc^2$ instead gives $\ell=2\hbar/(mc)$. The algebra decides neither.

% --- source: tex/s17b_rest#7 (Scale freedom)
\subsection{Scale freedom}
\label{sec:restscale}

The Clifford data $\eta$, $G$, $G\eta^{-1}$, $\eta G$, the operator $\cK$ of \cref{lem:Gclifford} and the structure constants of the complex cross
product $\beta_0$ of \cref{prop:lGhoms} are integer matrices, and the rescaling $(I_0,m)\mapsto(sI_0,tm)$, $s,t>0$, preserves admissibility and
multiplies $\varepsilon$ by $1/(st)$ at fixed initial rate $\kappa_0$, by $s/t$ at fixed $h(0)$ (the remark following \cref{prop:several}). Two
further scalings act on the time.

\begin{proposition}[Scalings]
\label{cl:scalings}
\begin{enumerate}[label=(\roman*),leftmargin=2.2em]
\item The interior law $\ddot X=-\kappa^2X$ is invariant under $X\mapsto\sigma X$, which multiplies $\tfrac m2\eta(\alpha_+X,\alpha_+X)$ by
$\sigma^2$, absorbed by $m\mapsto m/\sigma^2$; and under $(\tau,\kappa)\mapsto(s\tau,\kappa/s)$, which together with $m\mapsto sm$ preserves the
action.
\item The closure is invariant under $(\tau,I_0,m)\mapsto(s\tau,sI_0,sm)$, with $(\theta,h,p)$ unchanged as functions of the rescaled time;
$\kappa\mapsto\kappa/s$, while $m\kappa$, $\varepsilon$ and the action are invariant.
\end{enumerate}
\end{proposition}

\begin{proof}
(i) Linearity; $Y(\tau')=X(\tau'/s)$ solves $Y''=-(\kappa/s)^2Y$, the Lagrangian with $sm$ in place of $m$ is multiplied by $s^{-1}$, and
$d\tau'=s\,d\tau$. (ii) With $d/d\tau'=s^{-1}d/d\tau$, $I_0'=sI_0$ and $N'=N/s$, both equations \eqref{eq:LProtor} are reproduced, $H'=H/s$, and
$\kappa_0'=\kappa_0/s$ gives $\varepsilon'=\varepsilon$; $h$ is unchanged.
\end{proof}

The proposition is fixed by the algebra.

\begin{proposition}[Nothing in the algebra fixes $\ell$]
\label{cl:noscale}
\begin{enumerate}[label=(\roman*),leftmargin=2.2em]
\item Schur's lemma \cite{Schur} fixes only the form $c\,\id$ of the conversion operator (\cref{thm:onec}), and $c$ converts units of time into
units of length.
\item No monomial $c^{n_1}\hbar^{n_2}$ is a length, a mass or a time; with a mass $m$ adjoined, the only length is $\hbar/(mc)$, up to a number.
\item The interior systems with different $\kappa>0$ are carried into each other by a rescaling of $\tau$ that fixes every Clifford datum
(\cref{cl:scalings}(i)), so no condition formulated in the Clifford data distinguishes one value of $\kappa$, or of $\ell=c/\kappa$; in the
closure none fixes $\varepsilon$.
\end{enumerate}
\end{proposition}

\begin{proof}
(i) The commutant of the vector representation of $\so3$ in $M_3(\R)$ is $\R\,\id$. (ii) In the basis $(\mathsf M,\mathsf L,\mathsf T)$ the
dimensions $[c]=(0,1,-1)$ and $[\hbar]=(1,2,-1)$ span a plane containing none of $(1,0,0)$, $(0,1,0)$, $(0,0,1)$; adjoining $[m]=(1,0,0)$ gives
an invertible matrix, with $(0,1,0)=-[c]+[\hbar]-[m]$. (iii) The data are integer matrices, the rescaling acts on $\tau$ and $\kappa$ only, and
$\varepsilon$ is moved by the rescaling of $(I_0,m)$.
\end{proof}

Item (i) restates \cref{thm:onec}, fixed there, and item (ii) is inherited dimensional analysis, already implicit in \eqref{eq:kappa}, where the sum of $\partial_\lambda$ and $G$ requires a
scale; (iii) is fixed by the algebra. This is the one exact sense available here of the phrase that $m$ breaks the symmetry of $c^2$, the symmetric part of $E=mc^2$ in
the reading of \cref{sec:energy-readings} (compare \cref{sec:massverdict}): $c$ carries no
scale, the shift $\kappa G$ needs a length $\ell$ that the algebra cannot supply, and through $\ell=\hbar/(mc)$ that length would be a mass. What
it would break is the change of the unit of $\tau$ (\cref{cl:scalings}), not a symmetry of a vacuum; that is a reading.

% --- source: tex/s17b_rest#8 (The frame drifts)
\subsection{The frame drifts}
\label{sec:restdrift}

\Cref{ex:drift} computes the drift of the frame in closed form, $\Delta x=(8\mathrm K(-2)-4\mathrm E(-2))e_1$ per full turn of the rotor. After half a
turn $h$ has returned and $p=-p_0$, since $e^{\pi G}=-\id$; after a full turn both have returned. By \cref{thm:freedrift} every rotating solution
with $R_0>0$, that is with $\nu\neq0$ and a nonzero rest energy of the momentum field, has an unbounded frame, and by \cref{cl:nogen}(ii) the rest energy of the frame field grows with
it, as $\tfrac12k^2|\Delta x|^2$ after $k$ turns in the example. A constant nonzero rate and a bounded frame coexist, for a free rotor with
Lorentz-class inertia, exactly on null data (\cref{sec:closure}); at the equilibria of \cref{cl:flat} the rate is zero and the frame moves
uniformly; and with several constituents they coexist on the balance surface of \cref{prop:several}. Read in the terms of this section, and only as a reading, the differentiation
of the energy in motion unbinds the frame.

% --- source: tex/s17b_rest#9 (A quantum rotor, conditionally)
\subsection{A quantum rotor, conditionally}
\label{sec:restquantum}

\begin{proposition}[The half turn]
\label{cl:halfturn}
\begin{enumerate}[label=(\roman*),leftmargin=2.2em]
\item $e^{\pi G}=-\id$, $e^{2\pi G}=\id$, and for $p_0\neq0$ the leaf circle $\{e^{\theta G}p_0\}$ has period exactly $2\pi$.
\item For every symmetric $N$, $U(\theta+\pi)=U(\theta)$.
\item $\Gamma=e_1e_2\cdots e_6$ anticommutes with every $e_a$ and $\Gamma^2=\id$, so $\Ad_\Gamma=-\id$ on $V$ and $\Gamma$ implements $e^{\pi G}$ on
spinors. By \cref{cl:norel}(ii) the multiples of the half turn are the only elements of the rotor group that lift to $\mathrm{Pin}(3,3)$.
\end{enumerate}
\end{proposition}

\begin{proof}
(i) $(e^{\theta G}-\id)p_0=(\cos\theta-1)p_0+\sin\theta\,Gp_0$ with $p_0\perp Gp_0$ and $|Gp_0|=|p_0|$. (ii) $p(\theta+\pi)=-p(\theta)$ and $U$ is
quadratic. (iii) A product of all six anticommuting generators anticommutes with each of them, and
$(e_1\cdots e_6)^2=(-1)^{15}\prod_a\eta_{aa}=\id$.
\end{proof}

The proposition restates the remark on the lift of the rotor at the end of the three-constituent part of \cref{sec:closure} and makes its
ingredients explicit; it is fixed. The quantum statement it permits is conditional on an assumption.

\begin{assumption}[A quantised leaf]
\label{cl:quant}
The leaf $T^*S^1$ of \cref{prop:LP} (\cref{cl:frozen}(ii)) is quantised by wave functions with $\Phi(\theta+2\pi)=e^{2\pi i\delta}\Phi(\theta)$, $h=-i\hbar\,\partial_\theta$,
and the Casimirs of the leaf are treated as numbers.
\end{assumption}

\begin{proposition}[A conditional rate ladder]
\label{cl:ladder}
Under \cref{cl:quant}:
\begin{enumerate}[label=(\roman*),leftmargin=2.2em]
\item $h$ has spectrum $\hbar(\Z+\delta)$. On the leaves on which $U$ is constant ($R_0=0$, or duality-class inertia) the levels are
$E^{\rm rot}_n=\hbar^2(n+\delta)^2/(2I_0)$, plus the constant $\tfrac12p_0^{\!\top}Np_0$ in the duality class, the rates are
$\kappa_n=(n+\delta)\hbar/I_0$, and $E^{\rm rot}_{n+1}-E^{\rm rot}_n=\tfrac\hbar2(\kappa_n+\kappa_{n+1})$.
\item If $R_0>0$ the Hamiltonian is a Mathieu operator \cite{Mathieu,DLMF} whose potential couples $e^{i(n+\delta)\theta}$ only to
$e^{i(n\pm2+\delta)\theta}$, so for every $\delta$ the parity of $n$ is conserved; it labels the eigenvalues $\pm e^{i\pi\delta}$ of the half-turn
translation $\theta\mapsto\theta+\pi$. For $q_{\rm M}:=I_0R_0/(2m\hbar^2)\gg1$ the lowest gap in either parity sector is
$\hbar\Omega-\hbar^2/(2I_0)+O(q_{\rm M}^{-1/2}\hbar^2/I_0)$, and the two sectors are degenerate up to tunnelling.
\item If moreover $\delta=0$, and, as a separate assumption, the first level $h=\hbar$ is matched by $\hbar\kappa_1=mc^2$ with $\ell=\hbar/(mc)$, then
$I_0=\hbar^2/(mc^2)$, $M_{\rm rot}=m$ and $E^{\rm rot}_1=\tfrac12mc^2$.
\end{enumerate}
\end{proposition}

\begin{proof}
(i) Spectrum of $-i\hbar\partial_\theta$ on twisted periodic functions; in the duality class $U$ is the constant found above. (ii)
$\cos(2\theta+\varphi_0)$ shifts $n$ by $\pm2$, and $e^{i(n+\delta)(\theta+\pi)}=(-1)^ne^{i\pi\delta}e^{i(n+\delta)\theta}$. Dividing the Schr\"odinger
equation by $-\hbar^2/(2I_0)$ gives $\Phi''+\bigl(a_{\rm M}-2q_{\rm M}\cos(2\theta+\varphi_0)\bigr)\Phi=0$ with $a_{\rm M}=2I_0E/\hbar^2$, and
$a_{\rm M}\sim-2q_{\rm M}+2(2r+1)\sqrt{q_{\rm M}}-\tfrac18\bigl((2r+1)^2+1\bigr)$ as $q_{\rm M}\to\infty$ \cite[\S28.8]{DLMF}. The levels $r=0,1$ lie in
the same sector and differ by $4\sqrt{q_{\rm M}}-1$, while $\hbar\Omega=(\hbar^2/2I_0)\,4\sqrt{q_{\rm M}}$. (iii) $\hbar/I_0=\kappa_1=c/\ell=mc^2/\hbar$,
$M_{\rm rot}=I_0/\ell^2$ and $E^{\rm rot}_1=\hbar^2/(2I_0)$.
\end{proof}

Numerically, at $R_0/m=2000$ and $\hbar=I_0=1$ the lowest gap in each sector is $62.7415$ for $\delta=0$ and for $\delta=0.3$, against
$\hbar\Omega-\hbar^2/(2I_0)=62.7456$, and $62.7416$ with the next term of the expansion. The proposition is constructed. The quantisation of the
leaf exists: it is the one the orbit method assigns \cite{Kirillov04}, with $\delta$ a character of the stabiliser $2\pi\Z$ in the universal
cover of the rotor group. Nothing selects $\delta$, $I_0$ or the matchings, and whether the ladder is physical is open. Missing are a reason for
the assumption, that is a quantum theory of the momentum field, which the closure does not give and the split-form frame obstructs, since quantised
mode by mode with positive norm it has no ground state (\cref{mb:noground}). Missing too are
a value of $I_0$; a coupling to matter, since only the $\Z_2$ of the half turn reaches spinors, as $\Gamma$, and nothing ties the parity of $n$ to
chirality; and any derivation of the two matchings in (iii). The level $E^{\rm rot}_1=\tfrac12mc^2$ is $K_{\rm rot}=\tfrac12h\kappa$ at $h=\hbar$: the
factor $\tfrac12$ is built in, and $mc^2$ appears only as $\hbar$ times the rate put in, neither as a level nor as a transition energy.

\section{The identification tested}
\label{sec:esmap}

The preceding sections speak of a rotor, an inertia, a Lie--Poisson system and Euler--Arnold equations. This section states the
identification with a rigid body as a map, and tests whether the map preserves equations, brackets, conserved quantities and degrees of freedom.
The checks are rb2-1 to rb2-15.

\subsection{The map}
\label{sec:esmap:map}

\Cref{tab:esmap} sets the closure of \cref{sec:closure} beside the two rigid bodies of \cref{sec:rigidbody}. The closure is taken in the group
form of \cref{prop:groupred}: its configuration is $(\theta,x)$ in the group $S$ with product $(\theta,x)(\theta',x')=(\theta+\theta',x+e^{\theta G}x')$,
where $\theta$ is the rotor angle and $x\in V=\R^{3,3}$ the frame, with body form $X=e^{-\theta G}x$. Its Lagrangian is
$\tfrac12I_0\dot\theta^2+\tfrac m2\,\eta(u,u)$ in the body velocity $u=e^{-\theta G}\dot x$, or $\tfrac m2|u|^2$ for the positive form. The evolution
parameter is the interior parameter $\tau$, with the dimension of a time; this chapter does not relate it to the time read by a clock
(\cref{assum:clockred}).
With $X$ a length, the convention of \cref{cl:scalings}(i), the units are those of \cref{tab:esmap}, and both equations of \eqref{eq:LProtor}
are dimensionally consistent (rb2-12): $h$ is an angular momentum, $p$ a linear momentum and $I_0$ a moment of inertia.

\begin{table}[t]
\centering\small
\begin{tabular}{@{}>{\raggedright\arraybackslash}p{2.5cm}>{\raggedright\arraybackslash}p{3.3cm}>{\raggedright\arraybackslash}p{3.5cm}>{\raggedright\arraybackslash}p{4.7cm}@{}}
\toprule
 & rigid body & planar body & closure\\
\midrule
configuration & $R\in SO(3)$ & $(\theta,\mathbf x)\in SE(2)$ & $(\theta,x)\in S$: rotor angle and frame $x\in\R^{3,3}$\\
parameter & time $t$ & time $t$ & interior parameter $\tau$, $[\tau]=\mathsf T$\\
velocity & $\boldsymbol\Omega=(R^{-1}\dot R)^\vee$ & $(\dot\theta,\mathbf u)$, $\mathbf u=e^{-\theta J}\dot{\mathbf x}$ & $(\dot\theta,u)$, $u=e^{-\theta G}\dot x$\\
inertia & $\mathbb I>0$, $[\mathbb I]=\mathsf{ML}^2$ & $I>0$, $N^{-1}>0$ & $I_0>0$, $[I_0]=\mathsf{ML}^2$; frame $m\eta$ (split) or $m\,\id$ (positive), $[m]=\mathsf M$\\
action & $\int\tfrac12\boldsymbol\Omega\cdot\mathbb I\boldsymbol\Omega\,dt$ & $\int(\tfrac12I\dot\theta^2+\tfrac12\mathbf u^{\!\top}N^{-1}\mathbf u)\,dt$ & $\int(\tfrac12I_0\dot\theta^2+\tfrac m2\eta(u,u))\,d\tau$\\
momenta & $\mathbf M$, $[\mathbf M]=\mathsf{ML}^2\mathsf T^{-1}$ & $(h,p)$ & $(h,p)$, $[h]=\mathsf{ML}^2\mathsf T^{-1}$, $[p]=\mathsf{MLT}^{-1}$\\
equations & \eqref{eq:rb:euler} & \eqref{eq:rb:se2} & \eqref{eq:LProtor}\\
bracket & \eqref{eq:rb:lp} & \cref{sec:rb:se2} & \eqref{eq:restLPgen}\\
\bottomrule
\end{tabular}
\caption{The identification as a map: the closure of \cref{sec:closure} beside the free rigid body and the planar body of \cref{sec:rigidbody}.
Units with $X$ a length (\cref{cl:scalings}); checks rb2-12.}
\label{tab:esmap}
\end{table}

\subsection{What the map preserves}
\label{sec:esmap:preserve}

\begin{theorem}[The closure is a rigid body of the planar kind]
\label{thm:esmap}
\begin{enumerate}[label=(\roman*),leftmargin=2.2em]
\item \emph{Equations.} The bracket \eqref{eq:restLPgen} is the Lie--Poisson bracket of $\mathfrak s=\R\hat g\ltimes_GV$, and for every symmetric $N$
the equations of $H=h^2/(2I_0)+\tfrac12p^{\!\top}Np$ are \eqref{eq:LProtor}. On each pairing plane $(e_i,e_{i+3})$ the bracket is that of
$\mathfrak{se}(2)^*$, with $G$ acting as $J$. The group $S$ consists of the translations of $\R^6$ and the common rotation $e^{\theta G}$ of the three
pairing planes.
\item \emph{Brackets.} At $p\neq0$ the Poisson tensor of $\mathfrak s^*$ has rank two, with five independent Casimirs, the duality-invariant
functions of $p$ (\cref{cl:frozen}); those of $\so3^*$ and $\mathfrak{se}(2)^*$ have rank two with one Casimir. $\mathfrak s$ is solvable and $\so3$
is perfect, so every Lie algebra homomorphism $\so3\to\mathfrak s$ is zero, and every homomorphism $\mathfrak s\to\so3$ has an image of dimension at
most one.
\item \emph{Conserved quantities.} The energy $H$, the Casimirs, the spatial momentum $p_s=e^{\theta G}p$ and
\[
J_s=h+p_s^{\!\top}Gx
\]
are conserved for every symmetric $N$: the analogues of the energy, $|\mathbf M|^2$, and the linear and total angular momenta of a rigid body
that rotates and translates.
\item \emph{Degrees of freedom.} The configuration group, the dual Lie algebra, a generic leaf and the Casimirs have dimensions $(7,7,2,5)$ for the
closure, against $(3,3,2,1)$ for $SO(3)$ and for $SE(2)$. In all three the reduced motion has one degree of freedom; in the closure it is the
rotor angle, and the six frame coordinates are cyclic.
\item \emph{Inertia.} For duality-class inertia, $[N,G]=0$, the rotor and the frame decouple: $\dot h=0$ for all data. For Lorentz-class
inertia, $N=(\nu_1\eta^{-1}+\nu_2G\eta^{-1})/m$, they exchange angular momentum at the rate $\dot h=\Im(\nu w(p))/m$, and $N$ has the eigenvalues
$\pm|\nu|/m$, three of each: no nonzero Lorentz-class inertia is positive definite.
\item \emph{Energy.} With nonzero Lorentz-class inertia ($\nu\neq0$) $H$ is unbounded below on $\mathfrak s^*$, and on each leaf it is bounded below by
$-|\nu w(p_0)|/(2m)$.
\end{enumerate}
\end{theorem}

\begin{proof}
(i) For $\mu=(h,p)$ and $\nabla F=(F_h,\nabla_pF)\in\mathfrak s$, $[\nabla F_1,\nabla F_2]=F_{1,h}G\nabla_pF_2-F_{2,h}G\nabla_pF_1$, and
$-\langle\mu,[\nabla F_1,\nabla F_2]\rangle$ is \eqref{eq:restLPgen}; the equations are \cref{prop:LP}, for general $N$. On $(e_1,e_4)$, $G$ is
$\bigl(\begin{smallmatrix}0&-1\\1&0\end{smallmatrix}\bigr)$ and $\{h,p_1\}=-p_4$, $\{h,p_4\}=p_1$ (rb2-1, rb2-2, rb2-4). $G^{\!\top}=-G$, so
$e^{\theta G}\in SO(6)$ rotates each pairing plane by $\theta$. (ii) The rank is that of $\{h,p_j\}=(Gp)_j$, $\{p_i,p_j\}=0$; the five Casimirs are
\cref{cl:frozen}(ii) (rb2-3). $[\mathfrak s,\mathfrak s]=V$ and $[V,V]=0$. A homomorphism from the simple $\so3$ has kernel $0$ or $\so3$, and
its image cannot be a copy of $\so3$ inside the solvable $\mathfrak s$. The image of $\mathfrak s$ in $\so3$ is solvable, and $\so3$ has no
two-dimensional subalgebra, since for independent $\mathbf a,\mathbf b$ the vector $\mathbf a\times\mathbf b$ is nonzero and orthogonal to both (rb2-14).
(iii) $\dot p=-\dot\theta Gp$ gives $\tfrac{d}{d\tau}(e^{\theta G}p)=0$. With $\dot x=e^{\theta G}Np$ and $p_s^{\!\top}Ge^{\theta G}=p^{\!\top}G$, since $G$
commutes with $e^{\theta G}$ and $(e^{\theta G})^{\!\top}=e^{-\theta G}$, one has $\dot J_s=-p^{\!\top}GNp+p^{\!\top}GNp=0$ (rb2-5). (iv) Count from (ii)
(rb2-13). (v) If $[N,G]=0$ then $GN$ is antisymmetric and $p^{\!\top}GNp=0$. For the Lorentz class, $-p^{\!\top}GNp=\Im(\nu w(p))/m$ and
$\tfrac12p^{\!\top}Np=\Re(\nu w(p))/(2m)$, and the characteristic polynomial of $N$ is $(\lambda^2-|\nu|^2/m^2)^3$ (rb2-6, rb2-7). (vi) For $\nu\neq0$, along an
eigenvector of $-|\nu|/m$, $H\to-\infty$ (at $\nu=0$, $H=h^2/(2I_0)$). On a leaf $w(p)=e^{2i\theta}w(p_0)$, and $\Re(\nu w)\ge-|\nu w(p_0)|$ (rb2-8).
\end{proof}

The theorem is fixed. Items (i) and (iii) restate \cref{prop:LP,prop:groupred,thm:drift} in the language of \cref{sec:rigidbody}. The conserved $J_s$ is the momentum map
of semidirect-product mechanics \cite{HolmMarsdenRatiu98,MarsdenRatiu99}; it and items (ii), (iv), (v) and (vi) are computed here for this closure. A numerical integration of the full closure, rotor, momentum and
frame, with Lorentz-class inertia keeps $H$, $|p|^2$, $|w(p)|$, $p_s$ and $J_s$ to $2.4\times10^{-10}$ or better over $\tau\in[0,200]$, while $h$
varies over a range $0.737$: the rotor and the frame exchange angular momentum, and their total is conserved (rb2-15; \cref{tab:numerics}).

The verdict is therefore two-sided. As a structure the identification holds. The closure is a Lie--Poisson system of a left-invariant
Lagrangian on a group of rotations and translations, its reduced motion on every leaf is one rotation angle, and on each pairing plane its
bracket and its equations are those of the planar body of \cref{sec:rb:se2}. The pendulum of \cref{thm:pendulum} is the pendulum of
\eqref{eq:rb:se2} for an anisotropic translational inertia: on each pairing plane the Lorentz-class $N$ has the principal values $|\nu|/m$ and
$-|\nu|/m$, and $|w(p_0)|=|p_0|^2$ for $p_0$ in that plane. As an identification with a free rigid body it fails: the closure contains no rotation group $SO(3)$, and it maps to one at
most through its rotor. As an identification with a physical body it fails as well. A positive inertia that couples the rotor to the frame exists,
for instance $N=(2\,\id+\eta^{-1})/m$ (\cref{thm:inertiareq}(iv)), but it is not Lorentz invariant; the Lorentz-invariant inertia that produces
the pendulum is indefinite.

\subsection{Five distinctions}
\label{sec:esmap:distinctions}

Five pairs of notions must be kept apart by a rigid-body reading. In the arena they separate as follows.

\begin{enumerate}[label=(\arabic*),leftmargin=2.2em]
\item \emph{An internal rotation and the orientation of a body.} The rotor group $e^{\theta G}$ commutes with the whole of
$\mathfrak{gl}(3,\C)=\{A:[A,G]=0\}$: it is the centre $U(1)$ of $U(3)$ acting on $\C^3=\R^6$. It fixes no nonzero vector unless
$\theta\in2\pi\Z$, since $\det(e^{\theta G}-\id)=(2-2\cos\theta)^3$, and it is an isometry of $\eta$ only for $\theta\in\pi\Z$ (rb2-9, rb2-10;
\cref{cl:norel}(ii)). In the field-strength reading it is the duality rotation $z\mapsto e^{i\theta}z$ (\cref{thm:pendulum}). The orientation of a
body is an element of $SO(3)$ acting on the directions of space, and every such element fixes its axis. The rotor angle is a phase: it has no
axis, so it can point nowhere and precess about nothing.
\item \emph{A full dynamical rotor and a rate-clamped reduction.} The clamped interiors of \cref{sec:shifted,sec:neutral} are the reductions
$\theta=\kappa\tau$ of the closure (\cref{prop:groupred}). A clamped motion solves the closure exactly when $\dot h$ vanishes along it: always with
duality-class inertia, and with Lorentz-class inertia only when $\nu w(p_0)=0$, that is for a null momentum field (rb2-11; \cref{cor:constrate}).
Otherwise the clamped rate is the limit $\varepsilon\to0$ of a heavy rotor (\cref{cor:multiscale}), and the free rotor drifts the frame
(\cref{thm:freedrift}).
\item \emph{Bounded solutions and a Hamiltonian bounded below.} Every solution of the split-form interior is periodic (\cref{thm:neutral}), while
its energy has signature $(6,6)$ (\cref{r1:prop:shift}(iii)). In the closure the Hamiltonian is unbounded below on $\mathfrak s^*$ and bounded below on
each leaf (\cref{thm:esmap}(vi)). Boundedness of the orbits is a property of each leaf; a positive energy is not available.
\item \emph{A classical oscillation and a healthy quantum state.} The bound orbits, quantised mode by mode with positive norm, have no ground state
(\cref{mb:noground}).
\item \emph{An algebraic gap and a binding energy.} The gap $\hbar\Omega-\hbar^2/(2I_0)$ of \cref{cl:ladder}(ii) separates two eigenvalues of one
Mathieu operator on one leaf, under the assumption \cref{cl:quant}. The closure has no threshold of separated constituents against which it
would be a binding energy. On the tower the analogous comparison is \cref{e2:prop:threshold}: at rest the graded level $(m_1^2+m_2^2)^{1/2}$ lies below
the free threshold $m_1+m_2$, and the inequality is not a bound-state condition.
\end{enumerate}

Item (1) fixes how far a rigid-body vocabulary carries in this chapter. ``Rotor'' and ``inertia'' name the structure of \cref{thm:esmap}, and
``rotation'' names the internal phase of item (1); none of them refers to a body in space.

\section{From Rigid-Body Rotation to Planetary Orbital Dynamics}
\label{sec:orbital}

This section is a bounded extension. Rigid-body rotation and the orbital motion of a centre of mass are different systems. The first lives on
$SO(3)$, the second in $\R^3$, and in the Newtonian two-body problem they separate exactly for spherically symmetric bodies and are coupled by a
gravitational torque otherwise. A comparison with the arena therefore needs a map, which is stated first. The benchmark is Newtonian, and its
potential is assumed. The checks are rb3-1 to rb3-13, and the numerical runs are collected in \cref{tab:numerics}.

\subsection{Two maps}
\label{sec:orbital:maps}

\begin{definition}[The orbit map and the spin map]
\label{def:orbitmaps}
\begin{enumerate}[label=(\roman*),leftmargin=2.2em]
\item The \emph{orbit map} reads the interior configuration in one pairing plane, $z=x+iy$ for $X=x\,e_1+y\,e_4$, as the relative position of
two bodies in their orbital plane, and the interior parameter $\tau$ as their time.
\item The \emph{spin map} reads the rotor angle $\theta$ of the closure as the spin angle of a rigid body turning about an axis normal to its
orbital plane, measured from the line to the primary, and $I_0$ as its moment of inertia $C$ about that axis. The rotor momentum $h=I_0\dot\theta$
is then the spin angular momentum relative to the rotating line to the primary; the spin angular momentum itself is $h+Cn$ on a circular
orbit of mean motion $n$, and $h+C\dot f$ on an eccentric one, with $f$ the true anomaly.
\end{enumerate}
\end{definition}

Neither map is supplied by the arena. Both are readings, and the propositions below test what they carry.

\subsection{The Newtonian benchmark}
\label{sec:orbital:newton}

\begin{assumption}[Newtonian benchmark]
\label{assum:newton}
Two bodies of masses $M_1$ and $M_2$ interact through the potential $-GM_1M_2/r$. A non-spherical body is a rigid body with principal moments
$A<B<C$ whose axis of largest moment is normal to the orbital plane, and it moves on a Keplerian orbit about a point mass $M$ that it does not
perturb.
\end{assumption}

The assumption is the comparison system. It is not derived from the arena, which contains no gravitational potential: the curvature identities
of \cref{ch:ch11_curvature} come without a gravitational action. Its consequences are inherited
\cite{ArnoldMMCM,LandauLifshitz1,GoldreichPeale66,MurrayDermott99}.
\begin{itemize}
\item The centre of mass moves uniformly. The relative position obeys $\ddot{\mathbf r}=-\mu\,\mathbf r/r^3$ with $\mu=G(M_1+M_2)$, and conserves the
energy $E=\tfrac12|\dot{\mathbf r}|^2-\mu/r$, the angular momentum $\mathbf L=\mathbf r\times\dot{\mathbf r}$ and the Laplace--Runge--Lenz vector
$\mathbf A=\dot{\mathbf r}\times\mathbf L-\mu\,\mathbf r/r$. The bounded orbits are ellipses with a focus at the centre of force, and their period is
$2\pi\sqrt{a^3/\mu}$.
\item By Bertrand's theorem the inverse-square law and Hooke's law are the only central forces all of whose bounded orbits close
\cite{Bertrand1873,ArnoldMMCM}. Because $\mathbf A$ is conserved, a Kepler orbit does not precess. Apsidal precession needs a perturbation, such as
an oblate primary, a third body or relativity, and the benchmark contains none.
\item Two spherically symmetric bodies attract exactly as point masses and exert no torque on each other, so their spin and orbit separate. A
triaxial body feels the torque $-\tfrac32(GM/r^3)(B-A)\sin2(\theta-f)$, with $\theta$ its spin angle and $f$ the true anomaly. On a circular orbit of
radius $a$, the angle $\gamma=\theta-nt$ to the line to the primary obeys
\begin{equation}
\label{eq:spinorbit}
C\ddot\gamma=-\tfrac32\,\frac{GM}{a^3}\,(B-A)\sin2\gamma ,
\end{equation}
a pendulum in $2\gamma$ whose stable point $\gamma=0$, with the long axis towards the primary, is synchronous rotation, the Moon's state. On an
eccentric orbit the torque contains resonances at the spin rates $jn$, $j\in\tfrac12\Z$, of strengths $H_j(e)$, among them Mercury's $3{:}2$;
at larger asphericity the rotation can be chaotic \cite{GoldreichPeale66,MurrayDermott99,WisdomPealeMignard84}.
\item The planar model has no spin-axis precession, which needs a torque component in the orbital plane.
\end{itemize}

\subsection{The orbit map: Hooke, not Kepler}
\label{sec:orbital:hooke}

\begin{proposition}[Interior orbits under the orbit map]
\label{prop:hooke}
\begin{enumerate}[label=(\roman*),leftmargin=2.2em]
\item In a pairing plane the split-form interior is Hooke's law, $\ddot x=-\kappa^2x$ and $\ddot y=-\kappa^2y$. Its bounded orbits are ellipses
centred on the centre of force, $z(\tau+\pi/\kappa)=-z(\tau)$, all of period $2\pi/\kappa$ whatever their size. A Kepler ellipse has the centre of
force at a focus, and its period scales as $a^{3/2}$.
\item The interior conserves $\ell=x\dot y-y\dot x$, but as the Noether charge $m(\dot xy-\dot yx)$ of the boost of the mixed plane
$(e_1,e_4)$. The Euclidean rotation of that plane carries solutions to solutions but does not leave the Lagrangian invariant. The conserved
energy $\tfrac m2[(\dot x^2+\kappa^2x^2)-(\dot y^2+\kappa^2y^2)]$ is indefinite.
\item The arena's complex invariant $q(z)=z\cdot z$, which is $z^2$ in one plane, carries the orbit $z=c_-e^{-i\kappa\tau}+c_+e^{i\kappa\tau}$ onto an
ellipse with one focus at the origin, the other at $4c_-c_+$, semi-major axis $a_K=|c_-|^2+|c_+|^2$ and eccentricity $2|c_-||c_+|/a_K$. Away from
collisions, $z\neq0$, and with the new time $dt=\lambda|z|^2d\tau$, $\lambda>0$, the image $Q(t)=z(\tau)^2$ solves
\[
\frac{d^2Q}{dt^2}=-\mu_K\frac{Q}{|Q|^3},\qquad \mu_K=\frac{4\kappa^2a_K}{\lambda^2},\qquad E_K=-\frac{2\kappa^2}{\lambda^2} .
\]
For $|c_-|=|c_+|$ the orbit passes through the origin, and its image, of eccentricity one, is a regularised collision orbit.
\item The coupling $\mu_K$ depends on the orbit: at fixed $\lambda$ the images form a family of Kepler problems of one energy, one coupling per
orbit, not one Kepler problem with its family of orbits. Beyond one plane the image of an orbit is an ellipse with foci
$2c_-\cdot c_+\pm2\bigl((c_-\cdot c_-)(c_+\cdot c_+)\bigr)^{1/2}$. The origin is a focus only when $(c_-\cdot c_+)^2=(c_-\cdot c_-)(c_+\cdot c_+)$. For
$c_-=e_1$ and $c_+=e_2$ the image $q=2\cos2\kappa\tau$ passes through the origin while $|z|$ stays constant.
\end{enumerate}
\end{proposition}

\begin{proof}
(i) The Euler--Lagrange equations of $\tfrac m2[(\dot x^2-\dot y^2)-\kappa^2(x^2-y^2)]-m\kappa\,\tfrac{d}{d\tau}(xy)$, which is \eqref{eq:splitL}
on the plane by \cref{thm:EBsplit} (rb3-1). Periodicity and the half-period identity are read off the general solution, and Kepler's third law
scales as stated (rb3-2). (ii) Under $(\delta x,\delta y)=(y,x)$ the Lagrangian changes by $\tfrac{d}{d\tau}\bigl(-m\kappa(x^2+y^2)\bigr)$, and Noether's
theorem gives the charge. Under $(\delta x,\delta y)=(-y,x)$ the change has the nonzero Euler--Lagrange expression $2m(\ddot y+\kappa^2y)$ in
$x$, so it is not a total derivative (rb3-1). (iii) $z^2-4c_-c_+=(c_-e^{-i\kappa\tau}-c_+e^{i\kappa\tau})^2$, so
$|q|+|q-4c_-c_+|=|z|^2+|c_-e^{-i\kappa\tau}-c_+e^{i\kappa\tau}|^2=2(|c_-|^2+|c_+|^2)$ (rb3-3); this is Bohlin's observation \cite{Bohlin1911}. The equation
for $Q$ is the Levi-Civita transformation \cite{LeviCivita1920}. One checks it directly, using $|\dot z|^2+\kappa^2|z|^2=2\kappa^2a_K$ along the orbit
(rb3-4). (iv) $\mu_K\propto a_K$ by (iii); for instance $(c_-,c_+)=(1,\tfrac25)$ and $(1,\tfrac45)$ give $\mu_K=\tfrac{116}{25}$ and $\tfrac{164}{25}$ at
$\kappa=\lambda=1$ (rb3-5). With $\alpha_-=c_-\cdot c_-$, $\alpha_+=c_+\cdot c_+$, $\alpha_0=2c_-\cdot c_+$ and $u_-=\sqrt{\alpha_-}e^{-i\kappa\tau}$,
$u_+=\sqrt{\alpha_+}e^{i\kappa\tau}$, the image is $q=\alpha_0+u_-^2+u_+^2$ with $u_-u_+$ constant. So $q-\alpha_0\mp2u_-u_+=(u_-\mp u_+)^2$, and
$|q-F_+|+|q-F_-|=2(|\alpha_-|+|\alpha_+|)$ for $F_\pm=\alpha_0\pm2\sqrt{\alpha_-\alpha_+}$ (rb3-6).
\end{proof}

Items (i), (ii) and (iv) are fixed; (iii) is inherited and applied to the arena's own $q$ \cite{Bohlin1911,LeviCivita1920,ArnoldHuygens}.
In Bertrand's dichotomy the interior realises the Hooke family, not the Kepler family. The squaring map that carries one to the other is the
arena's invariant $q$, but the inverse-square law does not follow. The time change, with its free scale $\lambda$, is put in by hand. The coupling
$\mu_K$ is a property of each orbit, not of the force. And the map works plane by plane. In three dimensions the Kepler problem is regularised
through the Kustaanheimo--Stiefel map $\R^4\to\R^3$ under a bilinear constraint \cite{KustaanheimoStiefel65}; the arena's
$q\colon\C^3\to\C$ is not that map. The numerical test rb3-11 integrates an interior orbit with the new time and, independently, the
Kepler orbit of its image. They agree to $2.5\times10^{-9}$ over ten Kepler periods (\cref{tab:numerics}): the equivalence is exact, and it is
an equivalence of orbits, not of laws.

\subsection{The spin map: the duality pendulum as spin--orbit coupling}
\label{sec:orbital:spin}

\begin{proposition}[The duality pendulum under the spin map]
\label{prop:spinmap}
\begin{enumerate}[label=(\roman*),leftmargin=2.2em]
\item With $\theta=\gamma+\tfrac12(\pi-\varphi_0)$, $I_0=C$ and $R_0/m=\tfrac32(GM/a^3)(B-A)$, the duality pendulum of \cref{thm:pendulum} is the
circular-orbit equation \eqref{eq:spinorbit}. Its stable point $\psi=\pi$ is $\gamma=0$, synchronous rotation, and the frequency $\Omega$ of
\cref{cl:omega} is the libration frequency $\bigl(3(GM/a^3)(B-A)/C\bigr)^{1/2}$. The librating solutions are synchronous rotation with a libration,
and the rotating ones are non-synchronous spin with a modulated rate.
\item $R_0=0$ and duality-class inertia, under which the rotor decouples from the frame, correspond to $B=A$, under which the spin decouples from
the orbit. The closure conserves $J_s=h+p_s^{\!\top}Gx$ while $h$ varies, as the full two-body problem conserves the sum of the spin angular
momentum $h+Cn$ and the orbital angular momentum.
\item The correspondence holds only on the circular orbit. On an eccentric orbit the spin--orbit equation is not autonomous: at $e=0.2$ the
resonance strengths are $H_{1/2}(e)=-0.0995$, $H_1(e)=0.9013$, $H_{3/2}(e)=0.6397$ and $H_2(e)=0.3101$, against $H_j(0)=\delta_{j,1}$. The
duality pendulum is autonomous and has one family of equilibria.
\end{enumerate}
\end{proposition}

\begin{proof}
(i) Substitution, with $\sin(2\gamma+\pi)=-\sin2\gamma$ and $\Omega^2=2R_0/(mI_0)$ (rb3-7). (ii) The torque is proportional to $B-A$ (rb3-8); for the
closure see \cref{thm:esmap}(iii) and (v). (iii)
$H_j(e)=\tfrac1{2\pi}\int_0^{2\pi}(a/r)^3\cos(2f-2jnt)\,d(nt)$, with $t$ measured from pericentre, computed by quadrature; at $e=0.05$ it
agrees with the series $H_{1/2}(e)=-\tfrac e2+\tfrac{e^3}{16}$, $H_1(e)=1-\tfrac52e^2+\tfrac{13}{16}e^4$ and $H_{3/2}(e)=\tfrac72e-\tfrac{123}{16}e^3$
to $1.2\times10^{-6}$ (rb3-9).
\end{proof}

Items (i) and (ii) are fixed; the spin--orbit facts in (iii) are inherited \cite{GoldreichPeale66,MurrayDermott99}. The equivalence in (i) is one of
equations, and its coefficients come from different places: in the spin--orbit problem from the assumed potential and the shape of the body, in
the closure from the inertia and the rest energy of the momentum field, $R_0=2|\nu|E_0(\pi_0)$. The arena supplies the form of the pendulum, not the
orbit. In the closure the angular momentum $h$ is exchanged with the frame, whose motion is the drift of \cref{thm:freedrift}, a translation in
$\R^{3,3}$ and not an orbit. Two numerical runs make the boundary visible (\cref{tab:numerics}, \cref{fig:orbital}). On a circular orbit the
spin--orbit angle and the rotor of the closure, integrated independently from corresponding data, agree to $5.1\times10^{-10}$ rad over $50$ orbits
(rb3-12). On an orbit of eccentricity $0.2$, a body started at pericentre with spin rate $1.52\,n$ librates in the $3{:}2$ resonance for these
initial data, with mean spin rate $1.49988\,n$ over orbits $200$ to $1000$; the model has no tidal dissipation, so this is libration, not
capture. The duality pendulum with the same initial data keeps its own mean rate $1.51855\,n$, which is the
closed-form value for the circular orbit (rb3-13).

\subsection{Numerical examples}
\label{sec:orbital:numerics}

\Cref{tab:numerics} collects the numerical runs of this chapter: each tests an equivalence or its failure. Units are dimensionless. The orbital
runs take $GM=a=1$, so that the mean motion is $n=1$ and the orbital period $2\pi$; the closure takes $m=I_0=1$. The solver is the explicit
Runge--Kutta method DOP853 of order eight \cite{HairerNorsettWanner93}, as implemented in SciPy \cite{Virtanen2020}, with the tolerances shown.
The environment is recorded with the run outputs in \path{reproducibility/selected-results/v1.0.1/} (Python 3.13, NumPy 2.5, SciPy 1.18), and the
checks are listed in \cref{app:manifest}. The commands are
\texttt{python3} \path{verification/ch04_rigid_body/rb2_es_map.py} and \texttt{python3} \path{verification/ch04_rigid_body/rb3_orbital.py} \texttt{-{}-write-data}, which also writes
the data of \cref{fig:orbital}(b).

\begin{table}[t]
\centering\small
\begin{tabular}{@{}>{\raggedright\arraybackslash}p{1.2cm}>{\raggedright\arraybackslash}p{4.6cm}>{\raggedright\arraybackslash}p{1.6cm}>{\raggedright\arraybackslash}p{3.1cm}>{\raggedright\arraybackslash}p{3.5cm}@{}}
\toprule
check & system and initial data & tolerance & invariant drift & result\\
\midrule
rb2-15 & closure, $\nu=1+\tfrac i2$; $\theta=0$, $h=1.1$, $p=(1,0.3,-0.2,0.1,0.4,-0.5)$, $x=0$; $\tau\le200$ & $10^{-12}$ & $H$, $|p|^2$, $|w|$, $p_s$: $\le7.4\times10^{-11}$; $J_s$: $2.4\times10^{-10}$ & $h$ varies by $0.737$; rotor and frame exchange angular momentum\\
rb3-10 & Kepler, $\mu=a=1$, $e=0.5$, pericentre; $100$ periods & $10^{-12}$ & $E$: $4.5\times10^{-10}$; $|\mathbf L|$: $1.6\times10^{-10}$; $\mathbf A$: $1.9\times10^{-10}$ & return to pericentre within $3.3\times10^{-7}$; no apsidal precession\\
rb3-11 & interior orbit $(c_-,c_+)=(1,\tfrac25)$, $\kappa=\lambda=1$, with $dt=\lambda|z|^2d\tau$; Kepler orbit, $\mu_K=4.64$, of its image; $10$ periods & $10^{-12}$ & -- & images agree to $2.5\times10^{-9}$; $\tau=10\pi$ is $t=10\times2\pi(a_K^3/\mu_K)^{1/2}$\\
rb3-12 & spin--orbit, $e=0$, $(B-A)/C=0.01$, $\gamma(0)=0.5$, $\dot\gamma(0)=0$; closure, $N=0.015\,\eta^{-1}$, $p_0=e^1$; $50$ orbits & $10^{-11}$ & closure $H$: $1.1\times10^{-12}$ & angles agree to $5.1\times10^{-10}$ rad (libration)\\
rb3-13 & spin--orbit, $e=0.2$, $(B-A)/C=10^{-3}$, $\theta(0)=f(0)=0$, $\dot\theta(0)=1.52\,n$; closure from the same data; $1000$ orbits & $10^{-11}$ & orbital energy: $4.0\times10^{-8}$ & spin--orbit librates in $3{:}2$, mean rate $1.49988\,n$; closure $1.51855\,n$ (closed form $1.51855\,n$)\\
\bottomrule
\end{tabular}
\caption{Numerical runs of \cref{sec:esmap,sec:orbital}: DOP853, relative and absolute tolerance as shown, dimensionless units with $GM=a=1$ for the
orbits and $m=I_0=1$ for the closure. Drifts are maximal deviations over the run. Every entry is printed by the check named in the first column.}
\label{tab:numerics}
\end{table}

\begin{figure}[t]
\centering
\begin{tikzpicture}
\begin{axis}[name=left,width=6.6cm,axis equal image,xmin=-1.8,xmax=2.9,ymin=-1.3,ymax=1.3,
  axis lines=middle,xtick=\empty,ytick=\empty,xlabel={$\Re$},ylabel={$\Im$},
  title={\small (a) an interior orbit and its image under $q$}]
\addplot[tcol,thick,domain=0:360,samples=181] ({1.4*cos(x)},{-0.6*sin(x)});
\addplot[scol,thick,domain=0:360,samples=181] ({1.96*cos(x)^2-0.36*sin(x)^2},{-1.68*sin(x)*cos(x)});
\addplot[only marks,mark=*,mark size=1.6pt,black] coordinates {(0,0) (1.6,0)};
\node[tcol,anchor=south east,font=\small] at (axis cs:-1.25,0.45) {$z$};
\node[scol,anchor=south west,font=\small] at (axis cs:1.45,0.62) {$q=z^2$};
\node[anchor=north east,font=\scriptsize,inner sep=2pt] at (axis cs:1.62,0) {$4c_-c_+$};
\end{axis}
\begin{axis}[at={(left.north east)},anchor=north west,xshift=1.8cm,width=7.0cm,height=5.4cm,xmin=0,xmax=200,ymin=1.47,ymax=1.565,
  xlabel={orbits $t/2\pi$},ylabel={spin rate$/n$},
  title={\small (b) $e=0.2$: libration in $3{:}2$ and not},
  legend style={font=\scriptsize,at={(0.98,0.97)},anchor=north east,draw=none,fill=white,fill opacity=0.8,text opacity=1},legend cell align=left]
\addplot[scol,thick] table[x index=0,y index=1] {figures/ch04/spinrate_so.dat};
\addplot[tcol,thick] table[x index=0,y index=1] {figures/ch04/spinrate_es.dat};
\addplot[gray,dashed,domain=0:200] {1.5};
\legend{spin--orbit problem,duality pendulum,$3{:}2$ resonance}
\end{axis}
\end{tikzpicture}
\caption{(a) Analytical. The interior orbit $z=e^{-i\varphi}+\tfrac25e^{i\varphi}$ in one pairing plane (blue), an ellipse centred on the centre of
force, and its image $q=z^2$ (red), an ellipse with one focus at the origin and the other at $4c_-c_+=\tfrac85$ (\cref{prop:hooke}). (b) Numerical,
check rb3-13. The spin rate averaged over each orbit, $(\theta(t+2\pi)-\theta(t))/2\pi$, of a body with $(B-A)/C=10^{-3}$ on an orbit of
eccentricity $0.2$, started at pericentre at $1.52\,n$, over the first $200$ orbits: it librates about the $3{:}2$ resonance (red), between
$1.481\,n$ and $1.519\,n$, with a period of about $24$ orbits. The duality pendulum, which is the circular-orbit problem, keeps the rate
$1.5186\,n$ (blue). Data: \texttt{manuscript/figures/ch04/spinrate\_\{so,es\}.dat},
written by \texttt{python3 verification/\allowbreak ch04\_rigid\_body/\allowbreak rb3\_orbital.py} \texttt{-{}-write-data}; units $GM=a=1$, DOP853 with tolerance $10^{-11}$, $401$
samples.}
\label{fig:orbital}
\end{figure}

\subsection{What the extension shows}
\label{sec:orbital:verdict}

\begin{itemize}
\item \emph{Elliptical motion.} The interior has closed elliptical orbits, of Hooke's kind. Kepler ellipses appear only through the squaring map
$q$ and a time change put in by hand, with a coupling that varies from orbit to orbit. No inverse-square law is inferred, and no relativistic
correction; a familiar orbit can be drawn, which is not a derivation.
\item \emph{Conservation of energy and angular momentum.} The interior's energy is indefinite. Its angular momentum $\ell$ is a boost charge, not
the charge of a rotation of space.
\item \emph{Spin--orbit separation, torque and coupling.} Their form is present: duality-class inertia decouples rotor and frame, as a spherical
body decouples spin and orbit. Lorentz-class inertia couples them with the duality pendulum, as a triaxial body on a circular orbit is coupled.
The Lorentz-class inertia is indefinite, and the coupling constant is the rest energy of the momentum field, not a gravitational quantity.
\item \emph{Synchronous rotation.} The stable libration of the duality pendulum has the form of synchronous rotation. The resonances of
eccentric orbits, $3{:}2$ among them, have no counterpart.
\item \emph{Precession.} There is none. The rotor has no axis to precess about (\cref{sec:esmap:distinctions}, item (1)). The interior orbits
close, and the drift of the frame under the free rotor is a translation, not an apsidal precession.
\item \emph{The gravitational potential} is assumed throughout (\cref{assum:newton}). The arena derives none.
\end{itemize}
